\documentclass[11pt]{article}
\usepackage[margin=1in]{geometry}
\usepackage{amsmath,amssymb,bm}
\usepackage{physics}
\usepackage{hyperref}

\title{Notes on Entropy-Regularized Reinforcement Learning}
\author{}
\date{}

\begin{document}
\maketitle
\tableofcontents

\section{Equivalence Between Control and Inference}

\subsection*{Trajectory distributions}

The prior or uncontrolled trajectory distribution is
\begin{equation}
    p(\tau) = p(s_1)\prod_{t=1}^{N} p(s_{t+1}\mid s_t,a_t)\,\pi(a_t\mid s_t).
\end{equation}

A controlled trajectory distribution is
\begin{equation}
    p_c(\tau) = p_c(s_1)\prod_{t=1}^{N} p_c(s_{t+1}\mid s_t,a_t)\,\pi_c(a_t\mid s_t).
\end{equation}

The total reward of a trajectory is
\begin{equation}
    R_\tau = \sum_{t=1}^{N} r(s_t,a_t).
\end{equation}

In standard RL, the optimal policy is obtained by maximizing expected reward:
\begin{equation}
    \pi^*(a\mid s) = \arg\max_{\pi_c}\; \mathbb{E}_{p_c(\tau)}[R_\tau].
\end{equation}

In entropy-regularized RL, the objective becomes
\begin{equation}
    \mathbb{E}_{p_c(\tau)}[R_\tau] - \frac{1}{\beta}
    D_{\mathrm{KL}}\!\left(p_c(\tau)\middle\|p(\tau)\right),
\end{equation}
with
\begin{equation}
    D_{\mathrm{KL}}\!\left(p_c(\tau)\middle\|p(\tau)\right)
    = \sum_{\tau} p_c(\tau)\log\frac{p_c(\tau)}{p(\tau)}.
\end{equation}

\subsection*{Optimality variables and Bayes' rule}

The control-as-inference formulation introduces binary optimality variables \(O_t\) through
\begin{equation}
    p(O_t=1\mid s_t,a_t) = e^{\beta r(s_t,a_t)}.
\end{equation}

After shifting rewards so that \(r(s_t,a_t)\le 0\), the quantity
\(e^{\beta r(s_t,a_t)}\in (0,1]\), so it can be interpreted as the probability that step
\(t\) is optimal. Better state-action pairs correspond to larger probability of \(O_t=1\).

Define the event that the full trajectory is optimal by
\begin{equation}
    O_{1:N} = \bigcap_{t=1}^{N}(O_t=1).
\end{equation}

Assuming conditional independence of the optimality variables given the trajectory,
\begin{equation}
    p(O_{1:N}\mid \tau)
    = \prod_{t=1}^{N} p(O_t=1\mid s_t,a_t)
    = \prod_{t=1}^{N} e^{\beta r(s_t,a_t)}
    = e^{\beta R_\tau}.
\end{equation}

Bayes' rule then gives the posterior trajectory distribution
\begin{equation}
    p(\tau\mid O_{1:N})
    = \frac{p(O_{1:N}\mid \tau)\,p(\tau)}{p(O_{1:N})}
    = \frac{p(\tau)e^{\beta R_\tau}}{\sum_{\tau} p(\tau)e^{\beta R_\tau}}.
\end{equation}

If we define the energetic cost \(E_\tau=-R_\tau\), then
\begin{equation}
    p(\tau\mid O_{1:N})
    = \frac{p(\tau)e^{-\beta E_\tau}}{\sum_{\tau} p(\tau)e^{-\beta E_\tau}},
\end{equation}
which is Eq.~(6) in the paper.

\subsection*{Why this solves the KL-regularized control problem}

The optimality variable \(O_t\) is an auxiliary random variable. It does not add a new
physical degree of freedom to the environment. Its only purpose is to convert the control
problem into an inference problem.

Conditioning on \(O_{1:N}\) reweights the prior trajectory distribution toward larger-reward
trajectories:
\begin{equation}
    p(\tau\mid O_{1:N}) \propto p(\tau)e^{\beta R_\tau}.
\end{equation}

This posterior is also the optimizer of the variational problem
\begin{equation}
    \max_{q(\tau)}
    \left[
        \sum_{\tau} q(\tau)R_\tau
        - \frac{1}{\beta} D_{\mathrm{KL}}(q(\tau)\|p(\tau))
    \right].
\end{equation}

To see this, define
\begin{equation}
    q^*(\tau)=\frac{p(\tau)e^{\beta R_\tau}}{Z},
    \qquad
    Z=\sum_{\tau} p(\tau)e^{\beta R_\tau}.
\end{equation}

Then
\begin{equation}
    \log q^*(\tau)=\log p(\tau)+\beta R_\tau-\log Z,
\end{equation}
so
\begin{equation}
    \beta R_\tau-\log\frac{q(\tau)}{p(\tau)}
    = \log Z - \log\frac{q(\tau)}{q^*(\tau)}.
\end{equation}

Taking expectation with respect to any \(q(\tau)\), we obtain
\begin{equation}
    \beta \mathbb{E}_{q}[R_\tau] - D_{\mathrm{KL}}(q(\tau)\|p(\tau))
    = \log Z - D_{\mathrm{KL}}(q(\tau)\|q^*(\tau)).
\end{equation}

Therefore,
\begin{equation}
    \mathbb{E}_{q}[R_\tau]
    - \frac{1}{\beta}D_{\mathrm{KL}}(q(\tau)\|p(\tau))
    =
    \frac{1}{\beta}\log Z
    - \frac{1}{\beta}D_{\mathrm{KL}}(q(\tau)\|q^*(\tau)),
\end{equation}
and the objective is maximized exactly when \(q(\tau)=q^*(\tau)\).

Hence the posterior distribution obtained from Bayes' rule is exactly the same distribution
that solves the KL-regularized control problem.

\subsection*{Prior, posterior, and notation}

The prior trajectory distribution is
\begin{equation}
    p(\tau) = p(s_1)\prod_{t=1}^{N} p(s_{t+1}\mid s_t,a_t)\,\pi(a_t\mid s_t),
\end{equation}
with prior initial-state distribution \(p(s_1)\), prior dynamics \(p(s_{t+1}\mid s_t,a_t)\),
and prior policy \(\pi(a_t\mid s_t)\).

After conditioning on optimality, the posterior trajectory distribution is
\begin{equation}
    p(\tau\mid O_{1:N})
    = \frac{p(\tau)e^{\beta R_\tau}}{\sum_{\tau} p(\tau)e^{\beta R_\tau}}.
\end{equation}

Thus:
\begin{equation}
    p(\tau) \; \text{= prior/uncontrolled trajectory distribution},
    \qquad
    p(\tau\mid O_{1:N}) \; \text{= posterior/controlled trajectory distribution}.
\end{equation}

Before introducing \(O_{1:N}\), the paper writes the controlled distribution as \(p_c(\tau)\).
After introducing \(O_{1:N}\), the same object is written as the posterior
\(p(\tau\mid O_{1:N})\).

\subsection*{Constrained and unconstrained decompositions}

If only the policy is optimized, then
\begin{equation}
    p(\tau\mid O_{1:N})
    = p(s_1)\prod_{t=1}^{N} p(s_{t+1}\mid s_t,a_t)\,\pi(a_t\mid s_t,O_{1:N}).
\end{equation}

If the policy, dynamics, and initial-state distribution are all optimized, then
\begin{equation}
    p(\tau\mid O_{1:N})
    = p(s_1\mid O_{1:N})
    \prod_{t=1}^{N} p(s_{t+1}\mid s_t,a_t,O_{1:N})\,\pi(a_t\mid s_t,O_{1:N}).
\end{equation}

The difference is not the target object \(p(\tau\mid O_{1:N})\), but which factors are
allowed to change during optimization.

In the constrained setting,
\begin{equation}
    p(s_1\mid O_{1:N}) = p(s_1),
    \qquad
    p(s_{t+1}\mid s_t,a_t,O_{1:N}) = p(s_{t+1}\mid s_t,a_t),
\end{equation}
so only the posterior policy changes.

In the unconstrained setting, all factors are updated:
\begin{equation}
    p(s_1), \quad p(s_{t+1}\mid s_t,a_t), \quad \pi(a_t\mid s_t)
    \;\longrightarrow\;
    p(s_1\mid O_{1:N}), \quad p(s_{t+1}\mid s_t,a_t,O_{1:N}), \quad \pi(a_t\mid s_t,O_{1:N}).
\end{equation}

\subsection*{Why the optimal policy becomes stochastic}

The KL regularizer turns hard maximization into a soft reweighting. In the paper, the
optimal policy can be written as
\begin{equation}
    \pi^*(a\mid s)
    = \frac{\pi(a\mid s)\,u(s,a)}{\sum_{a'} \pi(a'\mid s)\,u(s,a')}.
\end{equation}

As long as multiple actions satisfy
\(
\pi(a\mid s)u(s,a)>0,
\)
they all retain positive probability. This is why KL-regularized optimization naturally
produces a stochastic policy.

\section{Introducing \texorpdfstring{$P$}{P}, Non-Absorption, and the Dominant-Eigenvalue Approximation}

\subsection*{The original transition matrix}

Let \(z=(s,a)\) and \(z'=(s',a')\) denote consecutive state-action pairs. The paper defines
\begin{equation}
    P_{ji}=p(z'=j\mid z=i)=p(s'\mid s,a)\pi(a'\mid s').
\end{equation}

This is a stochastic matrix because
\begin{equation}
    P_{ji}\ge 0,
    \qquad
    \sum_j P_{ji}=1
    \quad \text{for every } i.
\end{equation}

In the paper's convention, each column is a conditional probability distribution over next
state-action pairs.

\subsection*{The tilted matrix and loss of probability mass}

The tilted matrix is defined by
\begin{equation}
    \widetilde{P}_{ji} = P_{ji}e^{\beta r_i},
\end{equation}
where \(r_i=r(z=i)\).

Since the paper assumes \(r_i\le 0\), we have \(e^{\beta r_i}\le 1\), and therefore
\begin{equation}
    \sum_j \widetilde{P}_{ji}
    = e^{\beta r_i}\sum_j P_{ji}
    = e^{\beta r_i}\le 1.
\end{equation}

Thus \(\widetilde{P}\) is not stochastic in general. It is sub-stochastic: probability mass is
lost at each step, and the missing mass is
\begin{equation}
    1-\sum_j \widetilde{P}_{ji} = 1-e^{\beta r_i}.
\end{equation}

\subsection*{The extended matrix and non-absorption}

To recover a genuine stochastic process, the paper adds an absorbing state and defines the
extended matrix
\begin{equation}
    P =
    \begin{bmatrix}
        \widetilde{P} & 0 \\
        \delta & 1
    \end{bmatrix},
\end{equation}
where
\begin{equation}
    \delta_i = 1-\sum_j \widetilde{P}_{ji} = 1-e^{\beta r_i}.
\end{equation}

The missing probability mass is redirected into the absorbing state. If the process enters
that state, the trajectory is no longer regarded as optimal. Remaining in the original
state-action space means the trajectory has survived one more step without being absorbed.

Therefore, conditioning on trajectory optimality is equivalent to conditioning on
non-absorption for all \(N\) steps.

\subsection*{Meaning of \texorpdfstring{$[\widetilde{P}^{N}]_{ji}$}{[P-tilde\string^N]\_ji}}

For an ordinary stochastic matrix \(P\), \(P^N\) gives \(N\)-step transition probabilities.
For the sub-stochastic matrix \(\widetilde{P}\),
\begin{equation}
    [\widetilde{P}^{N}]_{ji}
\end{equation}
is the total weight of all \(N\)-step paths from \(i\) to \(j\) that survive without
absorption.

Summing over the final state gives the survival probability
\begin{equation}
    P(O_{1:N}\mid z_1=i)
    =
    \sum_j [\widetilde{P}^{N}]_{ji}.
\end{equation}

\subsection*{Primitive matrix}

The paper assumes that \(\widetilde{P}\) is primitive.

\paragraph{Definition.}
A square matrix \(A\in\mathbb{R}^{n\times n}\) with nonnegative entries is called primitive if
there exists an integer \(m\ge 1\) such that
\begin{equation}
    (A^m)_{ij}>0
    \qquad \text{for all } i,j.
\end{equation}

Equivalently, some power \(A^m\) is strictly positive.

This is \emph{not} the same definition as irreducible. Irreducibility only says that for each
pair \(i,j\), there exists at least one path length \(n_{ij}\) with
\((A^{n_{ij}})_{ij}>0\); the required length may depend on the pair. Primitivity is stronger in
form: it asks for one common exponent \(m\) such that every entry of \(A^m\) is strictly
positive at the same time. For a finite nonnegative matrix, the missing ingredient that turns
irreducibility into primitivity is aperiodicity. Precisely, for a finite nonnegative matrix,
primitivity is equivalent to:
\begin{itemize}
    \item irreducible: for any two indices \(i\) and \(j\), there exists some \(n\ge 1\) such
    that
    \begin{equation}
        (A^n)_{ij}>0;
    \end{equation}
    \item aperiodic: first define the period of state \(i\) by
    \begin{equation}
        d(i)=\gcd\{n\ge 1:(A^n)_{ii}>0\}.
    \end{equation}
    This means: look at all step numbers \(n\) for which the chain can return from \(i\) back
    to \(i\) in exactly \(n\) steps, and take their greatest common divisor. In an irreducible
    chain all states have the same period, so one can speak of \emph{the} period of the whole
    chain. The matrix is called aperiodic when this common period equals \(1\).
\end{itemize}

This definition is easier to understand by examples. If
\begin{equation}
    A=
    \begin{pmatrix}
        0 & 1\\
        1 & 0
    \end{pmatrix},
\end{equation}
then from state \(1\) one can return to state \(1\) only in \(2,4,6,\dots\) steps, so
\begin{equation}
    d(1)=\gcd\{2,4,6,\dots\}=2.
\end{equation}
The same is true for state \(2\). So this chain is irreducible but periodic, not aperiodic:
the motion is forced into an exact even-odd alternation.

By contrast, if there is any way to return in two different lengths that are not all multiples
of the same integer larger than \(1\), then the period becomes \(1\). For example, if a state
has a self-loop, then one can return in \(1\) step, and immediately
\begin{equation}
    d(i)=1.
\end{equation}
More generally, if returns are possible in \(2\) and \(3\) steps, then
\begin{equation}
    d(i)=\gcd(2,3)=1,
\end{equation}
so the chain is aperiodic. Intuitively, aperiodic means there is no rigid clock forcing all
returns to occur only at times divisible by some fixed integer \(d>1\).

The statement that all states in an irreducible chain have the same period is a standard but
important fact. Here is the reason. Suppose \(i\) and \(j\) communicate. Then there exist
integers \(m,n\ge 1\) such that
\begin{equation}
    (A^m)_{ji}>0,
    \qquad
    (A^n)_{ij}>0.
\end{equation}
Now take any return time \(r\) for state \(i\), meaning \((A^r)_{ii}>0\). By concatenating the
path
\begin{equation}
    j \to i \to i \to j,
\end{equation}
we obtain a return to \(j\) in exactly \(n+r+m\) steps, so
\begin{equation}
    (A^{n+r+m})_{jj}>0.
\end{equation}
Therefore every return time \(r\) of \(i\) produces a return time \(n+r+m\) of \(j\). This
implies that any common divisor of the return times of \(j\) must also divide differences of
the form
\begin{equation}
    (n+r_1+m)-(n+r_2+m)=r_1-r_2,
\end{equation}
where \(r_1,r_2\) are return times of \(i\). From this one gets
\begin{equation}
    d(j)\mid d(i).
\end{equation}
Repeating the same argument with \(i\) and \(j\) exchanged gives
\begin{equation}
    d(i)\mid d(j).
\end{equation}
Hence
\begin{equation}
    d(i)=d(j).
\end{equation}
Since irreducibility means every pair of states communicates, all states must share the same
period.

This assumption is exactly what ensures that the long-time behavior is governed by a single
dominant mode.

\subsection*{Perron-Frobenius theorem}

For a primitive nonnegative matrix \(A\in\mathbb{R}^{n\times n}\), a standard form of the
Perron-Frobenius theorem states:
\begin{quote}
There exists a real number \(\rho(A)>0\), called the Perron root, such that
\begin{equation}
    Av=\rho(A)v,
    \qquad
    u^\top A=\rho(A)u^\top
\end{equation}
for some vectors \(u,v\in\mathbb{R}^n\) with strictly positive components. Moreover,
\begin{equation}
    \rho(A)=\max\{|\lambda|:\lambda\in\sigma(A)\},
\end{equation}
\(\rho(A)\) is a simple eigenvalue, and every other eigenvalue \(\lambda\neq \rho(A)\)
satisfies
\begin{equation}
    |\lambda|<\rho(A).
\end{equation}
\end{quote}

Applying this to \(A=\widetilde{P}\), the paper obtains a unique dominant eigenvalue
\(\rho>0\), together with positive right and left eigenvectors \(v\) and \(u\):
\begin{equation}
    \widetilde{P}v=\rho v,
    \qquad
    u^\top\widetilde{P}=\rho u^\top.
\end{equation}

The chosen normalization is
\begin{equation}
    \sum_i v_i = 1,
    \qquad
    \sum_i u_i v_i = 1.
\end{equation}

\subsection*{Why the theorem works: proof sketch}

The theorem is a strict mathematical result. Its core logic is:
\begin{itemize}
    \item primitivity means some power of the matrix is strictly positive;
    \item repeated multiplication therefore pushes nonzero nonnegative vectors into the
    interior of the positive cone;
    \item this positivity and mixing select one preferred growth rate \(\rho\) and one
    preferred positive direction \(v\);
    \item applying the same argument to the transpose gives the positive left eigenvector
    \(u\);
    \item primitivity rules out competing modes on the spectral circle, so the Perron root is
    simple and strictly dominates the remaining eigenvalues.
\end{itemize}

\subsection*{Dominant-eigenvalue approximation}

Because every other eigenmode is suppressed relative to \(\rho^N\), the large-\(N\) behavior
of the matrix power is
\begin{equation}
    [\widetilde{P}^{N}]_{ji}
    \approx
    \rho^N u_i v_j.
\end{equation}

The paper writes the dominant eigenvalue as
\begin{equation}
    \rho=e^{-\beta\theta},
\end{equation}
so
\begin{equation}
    [\widetilde{P}^{N}]_{ji}
    \approx
    e^{-\beta\theta N}u_i v_j.
\end{equation}

Thus the long-time optimality and survival probabilities are controlled by a single scalar
\(\rho\) together with the two eigenvectors \(u\) and \(v\).

\subsection*{Detailed derivation of Eq.~(15) from Appendix B.1}

This is the step in the paper where the time-dependent posterior dynamics conditioned on
future optimality are converted into a time-independent ``driven'' Markov kernel. The clean
way to write it is as follows.

\paragraph{Step 1: notation.}
Let
\begin{equation}
    z_t=(s_t,a_t)\in\mathcal{Z}
\end{equation}
denote the state-action pair at time \(t\). The paper uses the index notation
\begin{equation}
    z_t=i,
    \qquad
    z_{t+1}=j.
\end{equation}
Write the future optimality event from time \(t\) onward as
\begin{equation}
    O_{t:T}=\{O_t=1,O_{t+1}=1,\dots,O_T=1\}.
\end{equation}
In the control-as-inference construction, \(O_t\) is an auxiliary binary variable with
distribution
\begin{equation}
    p(O_t=1\mid z_t=i)=e^{\beta r(i)},
\end{equation}
so even after conditioning on \(z_t=i\), the value of \(O_t\) is not deterministic. What is
true is weaker and more precise: once \(z_t=i\) is fixed, the random variable \(O_t\) carries
no additional information about future variables because it depends only on the current
state-action pair. Symbolically,
\begin{equation}
    O_t \perp (z_{t+1},\tau_{t+2:T},O_{t+1:T}) \mid z_t .
\end{equation}

\paragraph{A subtle but important point about the conditioning event.}
In the main text, the driven kernel is written with conditioning on the whole optimality
event \(O_{1:T}\). In Appendix B.1, however, the derivation is written using \(O_{t:T}\).
These are compatible, but one must prove the following conditioning reduction.

Write
\begin{equation}
    O_{1:T} = O_{1:t-1}\cap O_t \cap O_{t+1:T}.
\end{equation}
\paragraph{Proposition.}
Under the prior trajectory model,
\begin{equation}
    p(z_{t+1}=j\mid z_t=i,O_{1:T})
    =
    p(z_{t+1}=j\mid z_t=i,O_{t:T}).
    \label{eq:reduce-conditioning-window}
\end{equation}

\paragraph{What Markov does and does not imply.}
It is important not to misread \eqref{eq:reduce-conditioning-window} as a generic
``non-neighboring variables are independent'' statement. A first-order Markov property says
only that
\begin{equation}
    p(z_{t+1}\mid z_t,z_{t-1},\dots,z_1)=p(z_{t+1}\mid z_t),
\end{equation}
or equivalently,
\begin{equation}
    z_{t+1}\perp (z_{t-1},\dots,z_1)\mid z_t.
\end{equation}
This does \emph{not} imply that distant variables are unconditionally independent, and it
does \emph{not} imply that arbitrary conditioning events can be deleted. In the present proof,
the event \(O_{1:t-1}\) can be removed only because, after conditioning on \(z_t=i\), its
contribution enters as a common multiplicative factor in both numerator and denominator.
So Eq.~\eqref{eq:reduce-conditioning-window} is a special consequence of the particular
factorization of the augmented trajectory model, not a generic Markov-chain identity.

\paragraph{A useful corollary.}
If \(E\) is any event measurable with respect to the past history
\(\sigma(z_1,\dots,z_t)\), then the first-order Markov property does imply
\begin{equation}
    p(z_{t+1}\mid z_t,E)=p(z_{t+1}\mid z_t),
\end{equation}
whenever the conditional probabilities are well-defined. In words: once the current state
\(z_t\) is fixed, any extra condition built only from the past trajectory carries no further
information about the next step \(z_{t+1}\). However, this corollary does \emph{not} apply to
events involving future variables, such as \(O_{t+1:T}\), because those are not measurable
with respect to the past history.

\paragraph{Proof.}
Fix \(i,j\in\mathcal{Z}\). By definition,
\begin{equation}
    p(z_{t+1}=j\mid z_t=i,O_{1:T})
    =
    \frac{
        p(z_{t+1}=j,O_{1:t-1},O_t,O_{t+1:T}\mid z_t=i)
    }{
        p(O_{1:t-1},O_t,O_{t+1:T}\mid z_t=i)
    }.
    \label{eq:window-proof-start}
\end{equation}
We now use the factorization of the full joint law. Under the prior model,
\begin{equation}
    p(\tau_{1:T},O_{1:T})
    =
    p(z_1)\prod_{n=1}^{T-1}p(z_{n+1}\mid z_n)\prod_{n=1}^{T}p(O_n\mid z_n),
    \label{eq:full-joint-factorization}
\end{equation}
where \(p(O_n\mid z_n)\) depends only on \(z_n\). The issue is not that
\(O_{1:T}\) and \(O_{t:T}\) are the same event; they are not. The claim is only that they
induce the same conditional law for \(z_{t+1}\) once \(z_t=i\) is fixed.

To prove this without handwaving, let
\begin{equation}
    x=(z_1,\dots,z_{t-1}),
    \qquad
    y=(z_{t+2},\dots,z_T).
\end{equation}
Then
\begin{equation}
    p(z_{t+1}=j,O_{1:T},z_t=i)
    =
    \sum_{x,y} p(x,z_t=i,z_{t+1}=j,y,O_{1:T}).
\end{equation}
Using \eqref{eq:full-joint-factorization}, each summand factorizes as
\begin{align}
    &p(x,z_t=i,z_{t+1}=j,y,O_{1:T}) \nonumber\\
    &=
    \underbrace{
        p(z_1)\prod_{n=1}^{t-2}p(z_{n+1}\mid z_n)\,
        p(z_t=i\mid z_{t-1})
        \prod_{n=1}^{t-1}p(O_n=1\mid z_n)
    }_{\text{depends only on }x\text{ and }i}
    \nonumber\\
    &\quad\times
    p(O_t=1\mid z_t=i)\,
    p(z_{t+1}=j\mid z_t=i)\,
    \underbrace{
        \prod_{n=t+1}^{T-1}p(z_{n+1}\mid z_n)
        \prod_{n=t+1}^{T}p(O_n=1\mid z_n)
    }_{\text{depends only on }j\text{ and }y}.
    \label{eq:xy-factorization}
\end{align}
Now define
\begin{equation}
    A_i
    :=
    \sum_x
    p(z_1)\prod_{n=1}^{t-2}p(z_{n+1}\mid z_n)\,
    p(z_t=i\mid z_{t-1})
    \prod_{n=1}^{t-1}p(O_n=1\mid z_n),
\end{equation}
and
\begin{equation}
    H_j
    :=
    \sum_y
    \prod_{n=t+1}^{T-1}p(z_{n+1}\mid z_n)
    \prod_{n=t+1}^{T}p(O_n=1\mid z_n).
\end{equation}
Then \eqref{eq:xy-factorization} implies
\begin{equation}
    p(z_{t+1}=j,O_{1:T},z_t=i)
    =
    A_i\,
    p(O_t=1\mid z_t=i)\,
    p(z_{t+1}=j\mid z_t=i)\,
    H_j.
    \label{eq:joint-factor-numerator}
\end{equation}
Summing over \(j\) gives
\begin{equation}
    p(O_{1:T},z_t=i)
    =
    A_i\,
    p(O_t=1\mid z_t=i)\,
    \sum_{j'} p(z_{t+1}=j'\mid z_t=i)\,H_{j'}.
\end{equation}
Therefore
\begin{equation}
    p(z_{t+1}=j\mid z_t=i,O_{1:T})
    =
    \frac{
        p(z_{t+1}=j, O_{1:T}, z_t=i)
    }{
        p(O_{1:T}, z_t=i)
    }
    =
    \frac{
        p(z_{t+1}=j\mid z_t=i)\,H_j
    }{
        \sum_{j'} p(z_{t+1}=j'\mid z_t=i)\,H_{j'}
    }.
    \label{eq:reduced-fraction}
\end{equation}
The factors \(A_i\) and \(p(O_t=1\mid z_t=i)\) cancel completely.

On the other hand, using Bayes' rule directly for the shorter conditioning event gives
\begin{equation}
    p(z_{t+1}=j\mid z_t=i,O_{t:T})
    =
    \frac{
        p(z_{t+1}=j\mid z_t=i)\,
        p(O_t=1\mid z_t=i)\,
        H_j
    }{
        \sum_{j'} p(z_{t+1}=j'\mid z_t=i)\,
        p(O_t=1\mid z_t=i)\,
        H_{j'}
    }
    =
    \frac{
        p(z_{t+1}=j\mid z_t=i)\,H_j
    }{
        \sum_{j'} p(z_{t+1}=j'\mid z_t=i)\,H_{j'}
    }.
\end{equation}
Comparing with \eqref{eq:reduced-fraction} proves \eqref{eq:reduce-conditioning-window}.
\hfill\(\square\)

So Appendix B.1 is not changing the definition of the driven process; it is only using the
fact that, after conditioning on the current state-action pair \(z_t=i\), the past event
\(O_{1:t-1}\) and the current auxiliary variable \(O_t\) do not change the conditional law of
the next state-action pair beyond what is already contained in the future event \(O_{t:T}\).

\paragraph{Step 2: define the two kernels carefully.}
The paper defines the tilted kernel
\begin{equation}
    \widetilde{P}_{ji}
    =
    p(z_{t+1}=j\mid z_t=i)e^{\beta r(i)},
\end{equation}
which is time-independent, and the driven kernel
\begin{equation}
    [P_d^{(t,T)}]_{ji}
    :=
    p(z_{t+1}=j\mid z_t=i, O_{t:T}),
\end{equation}
which may depend on both \(t\) and \(T\) before the long-time limit is taken.

\paragraph{Step 3: factor the posterior transition in two different ways.}
Consider the joint conditional probability of the next state-action pair and the remaining
future path:
\begin{equation}
    p(z_{t+1}=j,\tau_{t+2:T}\mid z_t=i,O_{t:T}).
\end{equation}
Using the chain rule under the conditioned process gives
\begin{equation}
    p(z_{t+1}=j,\tau_{t+2:T}\mid z_t=i,O_{t:T})
    =
    p(\tau_{t+2:T}\mid z_{t+1}=j,z_t=i,O_{t:T})\,
    p(z_{t+1}=j\mid z_t=i,O_{t:T}).
\end{equation}
By the Markov property, once \(z_{t+1}=j\) is fixed, the future no longer depends on
\(z_t=i\). Also, after moving to time \(t+1\), the relevant optimality event becomes
\(O_{t+1:T}\). Hence
\begin{equation}
    p(\tau_{t+2:T}\mid z_{t+1}=j,z_t=i,O_{t:T})
    =
    p(\tau_{t+2:T}\mid z_{t+1}=j,O_{t+1:T}),
\end{equation}
so
\begin{equation}
    p(z_{t+1}=j,\tau_{t+2:T}\mid z_t=i,O_{t:T})
    =
    p(\tau_{t+2:T}\mid z_{t+1}=j,O_{t+1:T})\,[P_d^{(t,T)}]_{ji}.
    \label{eq:driven-factorization}
\end{equation}
This is Appendix B.1, Eq.~(B1), but with explicit time labels.

The logical content of this step is only the chain rule plus a conditional Markov property.
Nothing deeper is being used. More explicitly, the equality
\begin{equation}
    p(\tau_{t+2:T}\mid z_{t+1}=j,z_t=i,O_{t:T})
    =
    p(\tau_{t+2:T}\mid z_{t+1}=j,O_{t+1:T})
\end{equation}
holds because:
\begin{itemize}
    \item once \(z_{t+1}=j\) is fixed, the later path \(\tau_{t+2:T}\) no longer depends on
    \(z_t=i\);
    \item the variable \(O_t\) is not deterministic given \(z_t=i\), but it is conditionally
    independent of the future tail once \(z_t=i\) is fixed;
    \item the event \(O_{t+1}\) depends only on \(z_{t+1}=j\), so after conditioning on
    \(z_{t+1}=j\) it adds no further information about \(\tau_{t+2:T}\);
    \item the only remaining random constraints affecting the tail path are exactly the future
    events \(O_{t+1:T}\).
\end{itemize}
So the first factorization in Appendix B.1 is legitimate, but only after making this
conditioning reduction explicit.

\paragraph{Step 4: rewrite the same quantity using Bayes reweighting.}
From the posterior trajectory formula
\begin{equation}
    p(\tau_{t:T}\mid O_{t:T})
    \propto
    p(\tau_{t:T})e^{\beta\sum_{n=t}^{T}r(z_n)},
\end{equation}
conditioning additionally on \(z_t=i\) yields
\begin{equation}
    p(z_{t+1}=j,\tau_{t+2:T}\mid z_t=i,O_{t:T})
    =
    \frac{
        p(z_{t+1}=j,\tau_{t+2:T},O_{t:T}\mid z_t=i)
    }{
        p(O_{t:T}\mid z_t=i)
    }.
\end{equation}
Now use the fact that the immediate reward term \(e^{\beta r(i)}\) is attached to the
transition out of \(i\), and the remaining future event is \(O_{t+1:T}\). Then
\begin{equation}
    p(z_{t+1}=j,\tau_{t+2:T},O_{t:T}\mid z_t=i)
    =
    \widetilde{P}_{ji}\,
    p(\tau_{t+2:T},O_{t+1:T}\mid z_{t+1}=j).
\end{equation}
Therefore
\begin{equation}
    p(z_{t+1}=j,\tau_{t+2:T}\mid z_t=i,O_{t:T})
    =
    \frac{
        \widetilde{P}_{ji}\,
        p(\tau_{t+2:T},O_{t+1:T}\mid z_{t+1}=j)
    }{
        p(O_{t:T}\mid z_t=i)
    }.
    \label{eq:bayes-factorization}
\end{equation}
This is the content of Eq.~(B2).

\paragraph{Step 5: compare the two factorizations.}
Using
\begin{equation}
    p(\tau_{t+2:T}\mid z_{t+1}=j,O_{t+1:T})
    =
    \frac{
        p(\tau_{t+2:T},O_{t+1:T}\mid z_{t+1}=j)
    }{
        p(O_{t+1:T}\mid z_{t+1}=j)
    },
\end{equation}
and comparing \eqref{eq:driven-factorization} with \eqref{eq:bayes-factorization}, we obtain
\begin{equation}
    [P_d^{(t,T)}]_{ji}
    =
    \widetilde{P}_{ji}\,
    \frac{
        p(O_{t+1:T}\mid z_{t+1}=j)
    }{
        p(O_{t:T}\mid z_t=i)
    }.
    \label{eq:finite-time-driven}
\end{equation}
This is Appendix B.1, Eq.~(B3). Up to this point, nothing has been approximated.

It is useful to pause here and introduce a notation for the finite-time objects that appear
before any large-deviation limit is taken. Define the backward survival weight
\begin{equation}
    u_t(i) := p(O_{t:T}\mid z_t=i),
\end{equation}
and the forward weight
\begin{equation}
    v_t(i) := p(z_t=i,O_{1:t-1}).
\end{equation}
These are the natural finite-horizon analogues of the asymptotic left and right Perron
eigenvectors, but at this stage they are \emph{not} eigenvectors. They are simply conditional
and joint probabilities associated with a fixed finite horizon \(T\).

At finite horizon, the corresponding optimal policy is naturally time-dependent:
\begin{equation}
    \pi_t^*(a\mid s)
    =
    \frac{\pi(a\mid s)\,u_t(s,a)}
    {\sum_{\bar a}\pi(\bar a\mid s)\,u_t(s,\bar a)}.
\end{equation}
Only in the long-time limit, when \(u_t(s,a)\sim \rho^{\,T-t}u(s,a)\), does the common factor
\(\rho^{T-t}\) cancel so that one obtains the stationary form
\begin{equation}
    \pi^*(a\mid s)
    \propto
    \pi(a\mid s)\,u(s,a).
\end{equation}

With this notation, Eq.~\eqref{eq:finite-time-driven} can be rewritten as
\begin{equation}
    [P_d^{(t,T)}]_{ji}
    =
    \widetilde{P}_{ji}\frac{u_{t+1}(j)}{u_t(i)}.
\end{equation}
So the finite-time driven kernel is obtained by reweighting \(\widetilde P\) with the ratio of
future-optimality probabilities at two consecutive times. This is the exact finite-horizon
formula; no Perron-Frobenius approximation has entered yet.

The two weights satisfy simple recursions. The backward weight obeys
\begin{equation}
    u_t(i)=\sum_j \widetilde{P}_{ji}u_{t+1}(j),
    \qquad
    u_T(i)=p(O_T\mid z_T=i)=e^{\beta r(i)},
\end{equation}
or, depending on whether one chooses to attach the final reward to the last transition or to
the last state-action itself, an equivalent terminal normalization. In vector form this is the
backward recursion
\begin{equation}
    u_t = \widetilde P^\top u_{t+1}.
\end{equation}
Likewise, the forward weight satisfies
\begin{equation}
    v_{t+1}(j)=\sum_i \widetilde{P}_{ji}v_t(i),
\end{equation}
with \(v_1\) fixed by the initial state-action distribution. In vector form,
\begin{equation}
    v_{t+1}=\widetilde P\, v_t.
\end{equation}

Now the key probabilistic meaning of the product \(u_t(i)v_t(i)\) is:
\begin{equation}
    u_t(i)v_t(i)
    =
    p(z_t=i,O_{1:T}).
\end{equation}
Indeed,
\begin{equation}
    p(z_t=i,O_{1:T})
    =
    p(z_t=i,O_{1:t-1})\,p(O_{t:T}\mid z_t=i)
    =
    v_t(i)u_t(i).
\end{equation}
Therefore, after normalization over \(i\),
\begin{equation}
    p_t(i)
    :=
    \frac{u_t(i)v_t(i)}{\sum_k u_t(k)v_t(k)}
    =
    p(z_t=i\mid O_{1:T}).
\end{equation}
This is the exact time-\(t\) conditioned marginal distribution of the state-action pair in a
finite trajectory conditioned on global optimality.

This is precisely the object that later appears in the code for Figure 2. The arrays called
\texttt{u\_in\_time} and \texttt{v\_in\_time} are discrete implementations of these finite-time
backward and forward weights. Their entrywise product, after normalization, is not an
arbitrary heuristic quantity: it is the exact conditioned marginal
\begin{equation}
    p(z_t\mid O_{1:T})
\end{equation}
for the chosen finite horizon \(T\).

By contrast, the time-independent objects \(u(i)\) and \(v(i)\) that appear later are the
left and right Perron eigenvectors of \(\widetilde P\) in the long-time limit. Their product,
after the normalization
\begin{equation}
    \sum_i v_i=1,
    \qquad
    \sum_i u_i v_i=1,
\end{equation}
defines the bulk approximation
\begin{equation}
    p_{\mathrm{bulk}}(i):=u(i)v(i).
\end{equation}
So there is an essential distinction:
\begin{equation}
    p_t(i)=p(z_t=i\mid O_{1:T})
    \quad \text{is exact for finite } T,
\end{equation}
whereas
\begin{equation}
    p_{\mathrm{bulk}}(i)=u(i)v(i)
\end{equation}
is the asymptotic approximation valid in the interior bulk region of a long trajectory. The
paper's statement that
\begin{equation}
    \frac{u(s_t,a_t)v(s_t,a_t)}{p(s_t,a_t\mid O_{1:T})}\approx 1
\end{equation}
in the bulk region is precisely the claim that
\begin{equation}
    p_{\mathrm{bulk}}(z_t)\approx p(z_t\mid O_{1:T})
\end{equation}
when \(t\) is far from the two boundaries and \(T\) is large.

\paragraph{Step 6: express the survival probabilities by powers of \(\widetilde{P}\).}
Let \(N=T-t\). Then
\begin{equation}
    p(O_{t:T}\mid z_t=i)
    =
    \sum_{k\in\mathcal{Z}}[\widetilde{P}^{N}]_{ki},
\end{equation}
because we sum the weights of all \(N\)-step surviving paths starting from \(i\). Likewise,
\begin{equation}
    p(O_{t+1:T}\mid z_{t+1}=j)
    =
    \sum_{k\in\mathcal{Z}}[\widetilde{P}^{N-1}]_{kj}.
\end{equation}
Now assume \(\widetilde{P}\) is primitive, and let \(u,v\) be the positive left and right
Perron eigenvectors normalized by
\begin{equation}
    u^\top \widetilde{P}=\rho u^\top,
    \qquad
    \widetilde{P}v=\rho v,
    \qquad
    \sum_i v_i=1,
    \qquad
    \sum_i u_iv_i=1.
\end{equation}
Then for large \(N\),
\begin{equation}
    [\widetilde{P}^{N}]_{ki}
    \approx
    \rho^N v_k u_i,
\end{equation}
so summing over \(k\) gives
\begin{equation}
    p(O_{t:T}\mid z_t=i)
    \approx
    \rho^N u_i \sum_k v_k
    =
    \rho^N u_i,
\end{equation}
and similarly
\begin{equation}
    p(O_{t+1:T}\mid z_{t+1}=j)
    \approx
    \rho^{N-1}u_j.
\end{equation}

\paragraph{Step 7: take the long-time limit.}
Substituting these asymptotics into \eqref{eq:finite-time-driven} yields
\begin{equation}
    [P_d]_{ji}
    =
    \lim_{N\to\infty}[P_d^{(t,T)}]_{ji}
    =
    \widetilde{P}_{ji}\frac{\rho^{N-1}u_j}{\rho^N u_i}
    =
    \frac{\widetilde{P}_{ji}u_j}{\rho\,u_i}.
\end{equation}
Using the paper's notation
\begin{equation}
    \rho=e^{-\beta\theta},
\end{equation}
we obtain the main-text Eq.~(15):
\begin{equation}
    [P_d]_{ji}
    =
    \frac{\widetilde{P}_{ji}u_j}{e^{-\beta\theta}u_i}.
    \label{eq:doob-h-transform}
\end{equation}

\paragraph{Step 8: why this is called a Doob \(h\)-transform.}
If we define \(h(i)=u_i>0\), then \eqref{eq:doob-h-transform} can be written as
\begin{equation}
    [P_d]_{ji}
    =
    \frac{1}{\rho}\,
    \widetilde{P}_{ji}\,
    \frac{h(j)}{h(i)}.
\end{equation}
This is exactly the Doob \(h\)-transform of the sub-stochastic kernel \(\widetilde{P}\)
using the positive harmonic function \(h=u\). The role of the factor \(h(j)/h(i)\) is to
reweight transitions toward paths that are more compatible with long-time survival
(equivalently, long-time optimality).

\subsection*{Interpretation of the eigenvectors}

In the paper's later formulas:
\begin{itemize}
    \item \(u_i\) measures how favorable the current state-action pair \(i\) is for long-time
    survival or optimality;
    \item \(v_j\) describes the long-time distribution over surviving trajectories.
\end{itemize}

This is why the optimal policy is later expressed in terms of the left eigenvector \(u\).

\subsection*{Why this matters for the paper}

The point of the whole construction is that the large-\(N\) survival weights acquire a
universal form:
\begin{equation}
    [\widetilde{P}^{N}]_{ji}\sim \rho^N u_i v_j.
\end{equation}

Therefore
\begin{equation}
    P(O_{1:N}\mid z_1=i)
    =
    \sum_j [\widetilde{P}^{N}]_{ji}
    \approx
    \rho^N u_i,
\end{equation}
since the normalization \(\sum_j v_j=1\) is used.

This is the origin of the paper's statement that the long-time optimality probability decays
exponentially with \(N\), with a state-dependent prefactor \(u_i\).

\subsection*{References}

Useful references for the mathematical background are:
\begin{itemize}
    \item Levin and Peres, \emph{Markov Chains and Mixing Times}, cited as Ref.~[32] in the
    paper;
    \item van Doorn and Pollett, \emph{Quasi-stationary distributions for discrete-state
    models}, cited as Ref.~[34];
    \item Encyclopedia of Mathematics, ``Perron-Frobenius theorem'';
    \item Wolfram MathWorld, ``Perron-Frobenius Theorem''.
\end{itemize}

\section{Implementation Logic of the EntRegRL Project}

\subsection*{Overall map of the codebase}

The project is organized around four top-level Python files:
\begin{itemize}
    \item \texttt{frozen\_lake\_env.py}: defines the finite Markov decision process itself;
    \item \texttt{utils.py}: converts that MDP into sparse linear-algebra objects and solves the
    Perron-eigenvalue problem;
    \item \texttt{main.py}: runs the concrete experiments that produce the figures in the paper;
    \item \texttt{visualization.py}: turns arrays such as state distributions or policies into
    plots on top of the maze layout.
\end{itemize}

The cleanest way to read the code is not file by file in isolation, but as one pipeline:
\begin{equation}
    \text{map layout}
    \;\longrightarrow\;
    \text{finite MDP}
    \;\longrightarrow\;
    \text{sparse dynamics and rewards}
    \;\longrightarrow\;
    P,\widetilde{P}
    \;\longrightarrow\;
    (\rho,u,v)
    \;\longrightarrow\;
    \pi^*,\; p^*,\; \text{bulk distribution}
    \;\longrightarrow\;
    \text{figures}.
\end{equation}

So the four files are not four independent modules. They sit at four different conceptual
layers of the same derivation.

\subsection*{File 1: \texttt{frozen\_lake\_env.py} constructs the finite MDP}

This file answers the question: \emph{what is the uncontrolled process we are studying?}
Everything else in the project assumes that this object already exists.

The local class \texttt{DiscreteEnv} provides a minimal Gym-style discrete environment:
\begin{equation}
    \texttt{reset()} \to s_0,
    \qquad
    \texttt{step}(a_t) \to (s_{t+1}, r_t, done, truncated, info).
\end{equation}
This part is standard environment machinery: it stores the current state, samples from a
categorical transition distribution, and returns the realized next state and reward.

Since this terminology can be unfamiliar, it is worth unpacking it carefully.

Gym is a standard Python interface for reinforcement-learning environments. The basic idea is
that an environment is an object with a very small public API:
\begin{itemize}
    \item it can be reset to start a new episode;
    \item it can receive an action from the agent;
    \item it can return the next observation, reward, and termination flags.
\end{itemize}

In a \emph{discrete} environment, both the state space and the action space are finite. So
instead of continuous vectors, the code can label states and actions by integers:
\begin{equation}
    s \in \{0,1,\dots,n_S-1\},
    \qquad
    a \in \{0,1,\dots,n_A-1\}.
\end{equation}

The word \emph{minimal} means that the local class \texttt{DiscreteEnv} only implements the
small amount of machinery needed for this project. It is not a full-featured simulator. It
just provides the essentials.

The class stores several pieces of data:
\begin{itemize}
    \item \texttt{nS}: the number of states;
    \item \texttt{nA}: the number of actions;
    \item \texttt{P}: the transition dictionary;
    \item \texttt{isd}: the initial-state distribution;
    \item \texttt{s}: the current state of the environment.
\end{itemize}

The meaning of \texttt{isd} is especially important. It is a probability vector over initial
states:
\begin{equation}
    \texttt{isd}[s] = p(s_0 = s).
\end{equation}
So when the environment is reset, it does not necessarily always start from the same state; it
samples the initial state from this distribution.

Gym environments conventionally expose the types of actions and observations they support.
Here the code sets
\begin{equation}
    \texttt{action\_space} = \mathrm{Discrete}(n_A),
    \qquad
    \texttt{observation\_space} = \mathrm{Discrete}(n_S).
\end{equation}
This simply tells outside code that:
\begin{itemize}
    \item valid actions are integers from \(0\) to \(n_A-1\);
    \item observations are integers from \(0\) to \(n_S-1\).
\end{itemize}
In this project, the observation is just the state label itself, so there is no extra
observation-processing layer.

Calling \texttt{reset()} means: start a fresh trajectory. Concretely, the method samples
\begin{equation}
    s_0 \sim \texttt{isd},
\end{equation}
stores that sampled state as the new current state, clears the record of the last action, and
returns the state label.

So \texttt{reset()} is the operation that initializes one episode or one rollout.

Calling \texttt{step(a)} means: apply action \(a\) in the current state \(s\). The method then
looks up the list
\begin{equation}
    \texttt{P[s][a]} = [(p_1,s'_1,r_1,d_1),\dots,(p_m,s'_m,r_m,d_m)].
\end{equation}
This list represents all possible outcomes of taking action \(a\) in state \(s\).

The environment samples one of these outcomes according to the probabilities \(p_k\). If the
sampled outcome is \((p,s',r,d)\), then:
\begin{itemize}
    \item the current state is updated to \(s'\);
    \item the reward \(r\) is returned;
    \item the termination flag \(d\) is returned.
\end{itemize}

Here the symbol \(d_k\) means a boolean terminal flag attached to the \(k\)-th possible
transition outcome. If \(d_k=\texttt{True}\), that outcome ends the episode in the usual RL
sense; if \(d_k=\texttt{False}\), the trajectory continues. In this code there is also a
separate variable \texttt{truncated}, which is used when an outer \texttt{TimeLimit} wrapper
forces the rollout to stop after a fixed number of steps. So one can think of
\texttt{done}/\(d_k\) as ``the environment itself says stop'' and \texttt{truncated} as ``the
experiment script externally cuts the trajectory here''.

So \texttt{step(a)} is just a sampled realization of the MDP kernel
\begin{equation}
    p(s_{t+1},r_t,d_t \mid s_t,a_t).
\end{equation}
However, in the present implementation the reward is actually assigned before the next state is
used in any essential way. The code computes \texttt{rew} from the current cell type
\texttt{letter} and from whether the executed move is diagonal; it does not use
\texttt{newletter} to choose among different reward values. So for the environments used in
this project, the one-step reward is effectively a function of the current state-action pair,
\begin{equation}
    r(z)=r(s,a),
\end{equation}
rather than a genuinely next-state-dependent quantity \(r(s,a,s')\).

Even though \texttt{DiscreteEnv} is simple, it is sufficient because the rest of the project
only needs two perspectives on the same process:
\begin{itemize}
    \item a \emph{simulation view}, where trajectories are generated by repeated calls to
    \texttt{reset()} and \texttt{step()};
    \item a \emph{model view}, where the full transition dictionary \texttt{P} is read
    directly and converted into matrices.
\end{itemize}

This dual role is exactly why the class is useful. It is simultaneously:
\begin{itemize}
    \item an executable simulator of trajectories;
    \item an explicit tabulation of the finite MDP.
\end{itemize}

The real mathematical content appears in the subclass \texttt{ModifiedFrozenLake}. Its most
important output is the transition dictionary
\begin{equation}
    \texttt{P[s][a]} = \bigl[(p_1,s'_1,r_1,d_1),\dots,(p_m,s'_m,r_m,d_m)\bigr].
\end{equation}
This is a literal tabulation of the finite MDP. For each present state \(s\) and action
\(a\), the code stores every possible outcome:
\begin{itemize}
    \item transition probability \(p_k\),
    \item next state \(s'_k\),
    \item one-step reward \(r_k\),
    \item terminal flag \(d_k\).
\end{itemize}

Mathematically, this is the code representation of
\begin{equation}
    p(s_{t+1},r_t,\mathrm{done}_t \mid s_t,a_t).
\end{equation}
Later files never rebuild the environment from scratch; they simply read this explicit MDP.

The map is first stored as a character array \texttt{desc}. If the grid has shape
\((n_{\mathrm{row}},n_{\mathrm{col}})\), then the file defines
\begin{equation}
    n_S = n_{\mathrm{row}}n_{\mathrm{col}}.
\end{equation}
The helper
\begin{equation}
    \texttt{to\_s(row,col)} = row \cdot n_{\mathrm{col}} + col
\end{equation}
flattens each lattice site into one integer state label. This is important because every
later matrix in the project assumes that states are already encoded as integers
\(s \in \{0,\dots,n_S-1\}\).

The file does not treat the grid as a passive picture. Each symbol changes the transition law
or reward law:
\begin{itemize}
    \item \(S\): start state;
    \item \(F\): ordinary free cell;
    \item \(G\): goal;
    \item \(H\): hole;
    \item \(W\): wall;
    \item \(C\): candy, which increases reward;
    \item \(N\): nail, which decreases reward.
\end{itemize}
So the map is already a compact encoding of both geometry and cost landscape.

The constructor supports several action alphabets: \(2,3,4,5,8,9\) actions. The helper
\texttt{inc(row,col,action)} implements the kinematic rule
\begin{equation}
    (row,col,a) \mapsto (row',col').
\end{equation}
For example, with four actions the allowed moves are left, down, right, up; with five actions
there is also a stay action; with eight or nine actions diagonal moves are added.

This means that the file defines not just a reward function, but the underlying directed graph
on which trajectories move.

The helper \texttt{compute\_transition\_dynamics(...)} enumerates all outcomes for a chosen
intended action. If \texttt{slippery = 0}, then the executed action is deterministic. If
\texttt{slippery \(\neq 0\)}, the intended action is mixed with neighboring actions. Thus the
file directly builds the kernel
\begin{equation}
    p(s' \mid s,a)
\end{equation}
by explicit case analysis.

The reward in this environment is not a separate post-processing step. It is created at the
same time as the transition law. The code applies:
\begin{itemize}
    \item a step cost \(-\texttt{step\_penalization}\),
    \item an additional penalty on nails,
    \item an additional bonus on candy,
    \item a terminal bonus at goals,
    \item a terminal penalty at holes.
\end{itemize}
So each entry in \texttt{P[s][a]} already carries the one-step reward attached to the current
state-action pair. Even though the transition dictionary stores that reward alongside each
possible next state, in the present code the numerical reward value is the same across all
outcomes associated with a fixed current \(z=(s,a)\), except for differences caused by which
action is actually executed when the environment is slippery.

These two parameters are conceptually important for the paper:
\begin{itemize}
    \item \texttt{cyclic\_mode} makes hole states restart from the initial-state distribution
    instead of absorbing permanently;
    \item \texttt{never\_done} allows goal states to remain inside a continuing process rather
    than ending the episode.
\end{itemize}
This is not merely a coding convenience. The theory in the paper is fundamentally about
long-time trajectory ensembles and Perron eigenvectors of a continuing process, so the code
often prefers a non-terminating or effectively recycled dynamics.
\medskip
\texttt{frozen\_lake\_env.py} is the place where the abstract symbols
\(s,a,p(s'|s,a),r(s,a,s')\) become concrete finite arrays and dictionaries. It is the model
definition layer of the project.

\subsection*{File 2: \texttt{utils.py} converts the MDP into the paper's linear-algebra objects}

Once \texttt{env.P} is available, the next task is to rewrite the finite MDP in a form that
sparse linear algebra can handle efficiently. That is the job of \texttt{utils.py}. This file
contains the mathematical core of the implementation.

The function \texttt{get\_dynamics\_and\_rewards(env)} reads every entry of
\texttt{env.P[s][a]} and constructs two sparse objects.

First, it builds \texttt{dynamics} with shape
\begin{equation}
    (n_S,\; n_S n_A).
\end{equation}
The column index is
\begin{equation}
    i = s n_A + a,
\end{equation}
so each column corresponds to one state-action pair \(z=(s,a)\). The row index is the next
state \(s'\). Therefore
\begin{equation}
    \texttt{dynamics}[s',\, s n_A + a] = p(s' \mid s,a).
\end{equation}

This is a very important design choice. The code does not keep a 3-index tensor
\(p(s'|s,a)\). Instead it flattens the pair \((s,a)\) into one index, so that all later
operations can be expressed as sparse matrix products.

Second, it builds \texttt{rewards}. In principle, if an environment had a genuinely
next-state-dependent reward \(r(s,a,s')\), then one would need to average over \(s'\) to
produce a single state-action observable for each column. But in the present implementation the
reward is already effectively attached to the current state-action pair \(z=(s,a)\). So what
\texttt{get\_dynamics\_and\_rewards} stores is simply that state-action reward:
\begin{equation}
    \texttt{rewards}[s,a]=r(s,a).
\end{equation}
Because the code constructs the sparse matrix by summing all entries associated with the same
column, it may look formally like an expectation over next states, but in the environments used
here that sum is trivial: all nonzero contributions in a fixed \((s,a)\) column carry the same
reward value. Therefore no substantive averaging is taking place in the current project.

This state-action reward is the quantity that later enters the diagonal tilt
\(\mathrm{diag}(e^{\beta r_i})\).

The paper's main matrix \(P\) is not a state-to-state kernel. It is a state-action to
state-action kernel. The function \texttt{get\_mdp\_transition\_matrix(dynamics, policy)}
implements
\begin{equation}
    P_{ji} = p(z_{t+1}=j \mid z_t=i)
    = p(s'|s,a)\pi(a'|s').
\end{equation}

Operationally, this function does two things:
\begin{itemize}
    \item use \texttt{dynamics} to move from the current state-action pair \((s,a)\) to a next
    state \(s'\);
    \item use the policy to attach a next action \(a'\) to that next state.
\end{itemize}
So it turns the state-level environment kernel into a Markov chain on the enlarged
state-action space \(z=(s,a)\).

If \(D=n_S n_A\), then the resulting matrix has shape
\begin{equation}
    (D,D).
\end{equation}
This is the code object corresponding directly to the paper's \(P\).

These two functions implement the central mathematical statement of the paper:
\begin{equation}
    \widetilde{P} = P \, \mathrm{diag}(e^{\beta r_i}),
\end{equation}
then solve for its dominant eigenvalue and left/right Perron eigenvectors.

The code names the tilted matrix \texttt{M}. Its construction is:
\begin{enumerate}
    \item build \(P\) from \texttt{dynamics} and the prior policy;
    \item build the diagonal matrix
    \begin{equation}
        T = \mathrm{diag}(e^{\beta r_i});
    \end{equation}
    \item multiply
    \begin{equation}
        M = P T.
    \end{equation}
\end{enumerate}

This is exactly the paper's tilted operator, only under the variable name \texttt{M}.
The iterative updates are
\begin{equation}
    u_{k+1} \propto M^\top u_k,
    \qquad
    v_{k+1} \propto M v_k.
\end{equation}
So:
\begin{itemize}
    \item \texttt{u} approximates the left Perron eigenvector;
    \item \texttt{v} approximates the right Perron eigenvector;
    \item \texttt{l} approximates the dominant eigenvalue \(\rho\).
\end{itemize}

This matches the interpretation developed earlier in the note:
\begin{itemize}
    \item \(u\) contains the state-action dependent future survival or future optimality
    weight;
    \item \(v\) describes the asymptotic incoming distribution of surviving trajectories;
    \item \(\rho\) determines the exponential growth or decay rate of the tilted process.
\end{itemize}

The existence of both \texttt{solve\_unconstrained} and
\texttt{solve\_unconstrained\_v1} can look redundant, but conceptually they solve the same
problem. The difference is numerical stabilization:
\begin{itemize}
    \item \texttt{solve\_unconstrained} uses more aggressive rescaling heuristics, including an
    explicit \texttt{M\_scale}, to keep large-\(\beta\) runs stable;
    \item \texttt{solve\_unconstrained\_v1} is a cleaner power method with norm-based
    normalization.
\end{itemize}
So these are not different algorithms in the conceptual sense. They are two numerical
implementations of the same Perron-eigenproblem.

After obtaining \((l,u,v)\), the solver reconstructs the quantities that matter physically.

First, it computes the optimal policy by the Doob-transform-like formula
\begin{equation}
    \pi^*(a|s)
    =
    \frac{u(s,a)\pi(a|s)}
    {\sum_{a'} u(s,a')\pi(a'|s)}.
\end{equation}
This is precisely the policy reweighting predicted by the theory.

Second, it computes a state-level factor
\begin{equation}
    \chi(s) = \sum_a u(s,a)\pi(a|s),
\end{equation}
multiplies the state transition matrix by \(\chi\), and renormalizes columns. This produces
\texttt{optimal\_dynamics}, which is the code's realization of the driven process or optimal
dynamics.

Third, it normalizes \(u\) and \(v\) so that
\begin{equation}
    \sum_i v_i = 1,
    \qquad
    \sum_i u_i v_i = 1,
\end{equation}
and then forms the entrywise product \(u_i v_i\). Summing this over actions yields an
estimated bulk state distribution. This is the long-time occupancy measure in the tilted
ensemble.

\texttt{utils.py} also contains smaller but conceptually useful helpers:
\begin{itemize}
    \item \texttt{compute\_policy\_induced\_distribution}: propagates a policy forward for a
    fixed number of steps;
    \item \texttt{compute\_max\_kl\_divergence}: compares distributions;
    \item \texttt{test\_policy}: performs Monte Carlo rollouts in the actual environment.
\end{itemize}
These routines are used later in \texttt{main.py} for validation and figure generation.

\medskip
\texttt{utils.py} is the mathematical engine of the project. It takes the explicit finite MDP
and produces the paper's exact objects \(P\), \(\widetilde{P}\), \(\rho\), \(u\), \(v\),
\(\pi^*\), and the driven dynamics.

\subsection*{File 3: \texttt{main.py} is the experiment script that reproduces the paper}

\texttt{main.py} does not define the theory and does not define the environment. Its role is
to assemble those ingredients into concrete numerical experiments, each corresponding to one
or more figures in the paper.

So this file should be read as an executable appendix: each \texttt{make\_figure\_*} function
is a computational check of one claim in the paper.

Most figure functions begin with the same workflow:
\begin{enumerate}
    \item instantiate a \texttt{ModifiedFrozenLake} environment;
    \item optionally wrap it in \texttt{TimeLimit};
    \item call \texttt{get\_dynamics\_and\_rewards};
    \item choose a prior policy, usually uniform;
    \item either solve the Perron problem analytically or learn its eigenvectors from samples;
    \item convert the output into a plot.
\end{enumerate}
This repeated structure is useful because it makes clear that all figures are probing the same
underlying mathematical object from different angles.

\medskip
\texttt{make\_figure\_2} is the most direct numerical implementation of the paper's claim that
the interior of a long but finite trajectory is well approximated by the bulk distribution
built from the Perron eigenvectors. The code first fixes a concrete environment:
\begin{equation}
    \texttt{map\_name} = \texttt{'9x9zigzag'},
    \qquad
    N = 250,
    \qquad
    \beta = 20.
\end{equation}
The geometry of \texttt{9x9zigzag} is
\begin{verbatim}
FFFFFFFFF
FSFFFFFFF
WWWWWWFFF
FFFFFFFFF
FFFFFFFFF
FFFFFFFFF
FFFWWWWWW
FFFFFFFGF
FFFFFFFFF
\end{verbatim}
So this maze contains a start state \(S\), a goal state \(G\), many free cells \(F\), and two
horizontal wall barriers \(W\) that force a zigzag route. It contains no holes \(H\), no candy
\(C\), and no nails \(N\).

Then it builds \texttt{dynamics}, \texttt{rewards}, the uniform prior policy, and solves the
model-based Perron problem. From the solution it extracts the bulk state-action distribution
\begin{equation}
    p_{\mathrm{bulk}}(z) = u(z)v(z).
\end{equation}

For this specific map, the reward function is much simpler than the fully general rule defined
in \texttt{ModifiedFrozenLake}. In the code, Figure 2 uses
\begin{equation}
    \texttt{min\_reward}=-2,
    \qquad
    \texttt{step\_penalization}=1,
    \qquad
    \texttt{max\_reward}=0,
\end{equation}
while the default action set is the four cardinal moves and the default setting
\texttt{never\_done=True} is kept. Since the map contains only \(S\), \(F\), \(G\), and \(W\),
the one-step reward actually realized in this maze is:
\begin{itemize}
    \item \(r(z)=-1\) at every non-goal state-action pair, because each ordinary move carries a
    unit step cost;
    \item \(r(z)=0\) when the current state-action pair lies at the goal state, because the
    code sets the step penalty to zero once the process is already at \(G\);
    \item wall collisions do not introduce a special extra penalty: they simply keep the agent
    in place, but the step still costs \(-1\) because the current state is not the goal.
\end{itemize}
The parameter \texttt{min\_reward}=-2 does not play an active role in this particular maze,
because there are no hole states. Likewise, there is no candy or nail contribution. So for
Figure 2 the reward structure is essentially a shortest-path cost: every non-goal step costs
\(-1\), and being at the goal costs \(0\).

This matters because the tilted weight
\begin{equation}
    e^{\beta r(z)}
\end{equation}
is therefore just
\begin{equation}
    e^{-\beta}
\end{equation}
for ordinary non-goal moves and \(1\) at the goal. Since \(\beta=20\), non-goal steps are
penalized by an extremely small factor \(e^{-20}\). The conditioned ensemble thus strongly
favors trajectories that reach the goal region quickly and then exploit the zero-cost
recycling around the goal. That is the concrete reward-level reason why the bulk distribution
in Figure 2 concentrates on the effective optimal corridor of the maze.

The function then rebuilds the tilted matrix \(M=P\mathrm{diag}(e^{\beta r})\) explicitly and
starts a second calculation that is conceptually different from the eigenproblem. Instead of
computing only the asymptotic mode, it propagates finite-time forward and backward messages:
\begin{equation}
    v_{t+1} \propto M v_t,
    \qquad
    u_{t} \propto M^\top u_{t+1}.
\end{equation}
The code stores these iterates in \texttt{u\_in\_time} and \texttt{v\_in\_time}. At each time
slice \(t\), it multiplies them entrywise and renormalizes:
\begin{equation}
    p_t(z) \propto u_t(z)\,v_t(z).
\end{equation}
As explained earlier, this is exactly the finite-horizon conditioned marginal
\begin{equation}
    p_t(z)=p(z_t=z\mid O_{1:T}),
\end{equation}
not a heuristic proxy. In other words, the code is explicitly reconstructing the exact
distribution of the state-action pair at time \(t\), conditioned on the whole trajectory being
optimal. This \(p_t\) is the true time-dependent interior distribution for a finite horizon
\(N\), whereas \(p_{\mathrm{bulk}}\) is the large-\(N\) approximation.

Therefore the comparison in Figure 2 is exactly the same as the statement in the paper: the
product of the asymptotic Perron eigenvectors,
\begin{equation}
    u(z)v(z)=p_{\mathrm{bulk}}(z),
\end{equation}
is being compared with the exact finite-horizon conditioned marginal
\begin{equation}
    p_t(z)=p(z_t=z\mid O_{1:T}).
\end{equation}
The paper phrases the agreement by saying that the ratio
\begin{equation}
    \frac{u(z_t)v(z_t)}{p(z_t\mid O_{1:T})}
\end{equation}
is close to \(1\) in the bulk region. The code expresses the same idea through the KL
divergence between these two distributions as a function of \(t\).

The final numerical comparison is a KL divergence at each time \(t\). The code uses the bulk
distribution as \(p\) and the exact finite-time distribution as \(q_t\), computing
\begin{equation}
    D_{\mathrm{KL}}(t)
    =
    \sum_z p_{\mathrm{bulk}}(z)\log \frac{p_{\mathrm{bulk}}(z)}{p_t(z)}.
\end{equation}
So the plotted curve is not reward, policy, or value. It is a direct measure of how close the
bulk approximation is to the exact conditioned finite-time distribution as one moves from the
left boundary of the trajectory to the right boundary. The logic is that the center of the
trajectory should be well described by the Perron bulk, while the two ends should show
boundary effects.

\medskip
\texttt{make\_figure\_3} is the function where the paper's thermodynamic interpretation is made
fully explicit in a concrete finite system. It again starts from a fixed model, namely the same
\texttt{9x9zigzag} map,
and scans over many values of \(\beta\). For each \(\beta\), it solves the Perron problem and
stores the solution. Then it builds the state-action bulk distribution
\begin{equation}
    p_\beta(z)=u_\beta(z)v_\beta(z)
\end{equation}
and computes the quantities
\begin{equation}
    F,\qquad E,\qquad S,\qquad \theta=-\log\rho.
\end{equation}
More concretely, the code defines
\begin{equation}
    E(\beta) = -\sum_z p_\beta(z)\,r(z),
\end{equation}
which is the mean cost per step, and
\begin{equation}
    F(\beta) = -\frac{\log \rho(\beta)}{\beta},
\end{equation}
which plays the role of a free energy per step. It also computes an entropy-like quantity by
combining the optimal policy and the prior policy:
\begin{equation}
    S_1 = \sum_z p_\beta(z)\bigl[-\log \pi^*(a|s)\bigr],
    \qquad
    S_2 = \sum_z p_\beta(z)\log \pi(a|s),
\end{equation}
so that \(-(S_1+S_2)\) is the relative-entropy contribution plotted in the second panel.

The function makes two different kinds of figures. The top figure is a family of state
distributions for several representative values of \(\beta\). To obtain them, the code sums
\(p_\beta(z)\) over actions and uses \texttt{plot\_dist} to display the resulting state
occupancies on the maze. This gives an immediate visual answer to the question: as \(\beta\)
changes, where does the controlled ensemble spend its time?

The bottom figure is a thermodynamic summary. The code puts \(F(\beta)\), \(E(\beta)\), and the
entropy term on one set of axes, and then adds Monte Carlo points labelled ``Simulated
Energy''. The label can be slightly confusing, because what is actually simulated is the mean
reward per step and then converted into a mean cost per step so that it can be compared with
\(E(\beta)\).

More precisely, the helper \texttt{test\_policy(env,policy)} returns the total realized reward
of one rollout,
\begin{equation}
    R_{\mathrm{rollout}}=\sum_{t=1}^{N} r_t.
\end{equation}
In \texttt{make\_figure\_3}, the worker function records
\begin{equation}
    -\frac{R_{\mathrm{rollout}}}{N},
\end{equation}
where \(N\) is the trajectory length used in the \texttt{TimeLimit} wrapper. So each Monte
Carlo sample is a \emph{sample mean cost per step}. Averaging many such samples yields
\begin{equation}
    E_{\mathrm{sim}}(\beta)
    \approx
    -\frac{1}{N}\,\mathbb{E}_{\pi_\beta^*}[R_{\mathrm{rollout}}],
\end{equation}
which is exactly the quantity plotted as ``Simulated Energy''.

So these points are not produced by a separate RL algorithm. The code first computes the
analytical policy \(\pi^*_\beta\) from the Perron problem, then rolls out that policy in the
actual environment, measures its realized mean reward per step, flips the sign to convert
reward into cost, and compares the result with the analytical prediction
\begin{equation}
    E(\beta)=-\sum_z p_\beta(z)r(z).
\end{equation}
Figure 3 therefore simultaneously shows the thermodynamic quantities extracted from the Perron
problem and verifies, by direct simulation, that the analytically reconstructed policy has the
predicted average cost.

It is also important not to over-interpret the word ``thermodynamic'' here. All of these curves
are computed on one fixed finite \(9\times 9\) maze. So this is not a thermodynamic limit in
the spatial sense \(L\to\infty\). The relevant asymptotic idea in this figure is the
long-trajectory or large-deviation limit in time, not an infinite-volume limit in lattice size.

\medskip
\texttt{solve\_dp} and \texttt{make\_figure\_5} form the dynamic-programming comparison. Here
the environment is changed to \texttt{10x10empty}, and the goal is no longer to visualize a
policy or a state distribution. Instead, the code wants to compare two different ways of
computing soft \(Q\)-values.

The helper \texttt{solve\_dp(beta,max\_steps,gamma)} performs a finite-horizon soft
dynamic-programming recursion. Starting from zero \(Q\)-values and zero soft values, it
iterates
\begin{equation}
    Q(s,a) = r(s,a) + \gamma \sum_{s'} p(s'|s,a) V(s')
\end{equation}
followed by the soft backup
\begin{equation}
    V(s)=\frac{1}{\beta}\log \sum_a \pi(a|s)e^{\beta Q(s,a)},
\end{equation}
implemented with the usual numerical stabilization via a local shift \(\delta\).

The actual figure is produced by \texttt{make\_figure\_5}. This function computes the entire
finite-horizon sequence \(Q_t\) and \(V_t\) for all horizons up to \(N\). In parallel, it also
solves the large-deviation Perron problem and constructs the approximation
\begin{equation}
    Q^{\mathrm{LD}}_t(z)
    =
    \frac{\log \rho \, t + \log u(z)}{\beta},
\end{equation}
then divides by \(t\) to compare per-step quantities.

The resulting figure has several panels. One panel plots the root-mean-square error between the
finite-horizon soft \(Q\)-values and the large-deviation approximation as a function of horizon
length. Two other panels make direct scatter plots of
\begin{equation}
    Q^{\mathrm{DP}}_t(z)/t
    \quad \text{versus} \quad
    Q^{\mathrm{LD}}_t(z)/t
\end{equation}
for a moderately short horizon and for a long horizon. The dashed diagonal is the line of
perfect agreement. So Figure 5 is the place where the code demonstrates that the eigenvalue
formulation is not an unrelated alternative to soft value functions, but their long-time
asymptotic limit.

Here the role of \(N\) is completely explicit. First, the environment is wrapped as
\texttt{TimeLimit(env, N)}, so \(N\) defines the finite horizon of the control problem being
solved. Second, the code stores the whole family
\(\{Q_t^{\mathrm{DP}}\}_{t=1}^N\) and \(\{V_t^{\mathrm{DP}}\}_{t=1}^N\), so the top-right
panel is literally a plot of the error versus trajectory length \(t\le N\). Third, the two
bottom scatter plots are obtained by choosing two specific horizons, \(t=20\) and \(t=290\),
and comparing all state-action values at those two lengths. So in Figure 5, \(N\) is part of
the theoretical object itself: the figure compares finite-horizon dynamic programming with the
large-\(N\) Perron approximation.

\medskip
\texttt{make\_figure\_6} is the simplest figure in conceptual terms, but it is visually very
useful. The function loops over three different maps:
\begin{equation}
    \texttt{'7x7holes'},\qquad
    \texttt{'8x8zigzag'},\qquad
    \texttt{'9x9ridgex4'}.
\end{equation}
For each map it solves the model-based Perron problem at fixed \(\beta\) and extracts the
optimal policy \(\pi^*(a|s)\). It then calls \texttt{plot\_dist(env.desc, optimal\_policy)},
which interprets the \((n_S,n_A)\) table as arrows on the maze. No additional statistics are
computed here. The whole point of the figure is to show, directly and qualitatively, how the
left eigenvector reweights actions in different regions of the maze. In other words, Figure 6
is where the abstract Doob-transform formula turns into an actual navigation strategy visible
by eye.

\medskip
\texttt{make\_figure\_7} is the first genuinely sample-based learning routine. It uses the map
\texttt{'7x7holes'} and an episode length \(N=1000\). The function still calls
\texttt{get\_dynamics\_and\_rewards} once, but this is only to obtain dimensions and maintain
consistency of notation; the actual learning loop does not build or solve \(\widetilde P\).
Instead it begins by collecting a replay buffer of SARSA-style tuples
\begin{enumerate}
    \item sample trajectories under the uniform prior policy;
    \item store every encountered tuple \((s,a,r,s',a')\) in \texttt{replay\_memory};
    \item for each \(\beta\) in \([1,5,10,50,100]\), repeatedly sweep through that buffer;
    \item update a table \(u(s,a)\) and an eigenvalue estimate \(l\);
    \item convert the current \(u\) into a policy and test it by rollouts.
\end{enumerate}

The core update is
\begin{equation}
    u(s,a)
    \leftarrow
    (1-\alpha)u(s,a) + \alpha \frac{e^{\beta r}}{l}u(s',a').
\end{equation}
At the same time, the code updates
\begin{equation}
    l
    \leftarrow
    (1-\alpha_\theta)l + \alpha_\theta e^{\beta r}\frac{u(s',a')}{u(s,a)},
\end{equation}
with \(\alpha_\theta = \alpha /(n_S n_A)\). These are not arbitrary TD rules. They are the
sampled analogues of the left-eigenvector equation
\begin{equation}
    \rho\,u(s,a)
    =
    e^{\beta r(s,a)}
    \mathbb{E}[u(s',a') \mid s,a].
\end{equation}
Every fixed number of training steps, the code turns the learned \(u\)-table into a policy by
\begin{equation}
    \pi_t(a|s)\propto u_t(s,a)\pi(a|s),
\end{equation}
runs \texttt{test\_policy} ten times, and records the empirical mean reward per step. The
final plot is therefore a family of learning curves: horizontal axis is training step,
vertical axis is mean per-step reward, and each curve corresponds to one value of \(\beta\).
So Figure 7 is not showing convergence of the exact solver; it is showing whether a
sample-based \(u\)-learning rule can progressively improve control performance.

In Figure 7, \(N\) enters differently from Figure 5. It is no longer the horizontal axis of a
finite-horizon comparison. Instead, it defines the episode length used to collect data. Because
the environment is wrapped in \texttt{TimeLimit(env, N)}, each episode contributes at most
\(N\) SARSA tuples \((s,a,r,s',a')\) to the replay buffer. Hence the total amount of sampled
experience grows roughly like \(N\times n_{\mathrm{episodes}}\). The same \(N\) also sets the
evaluation and rescaling periods through
\texttt{evaluation\_period = N * 100} and \texttt{rescale\_period = N * 100}, and the measured
return is divided by \(N\) to produce a per-step reward. So in Figure 7, \(N\) is part of the
sampling and evaluation protocol rather than part of the Perron fixed-point equation itself.

\medskip
\texttt{make\_figure\_8} extends the previous idea by trying to learn both the left and the
right Perron eigenvectors from trajectories. It again uses \texttt{'7x7holes'} and long
episodes, but now launches several independent replicas in parallel on CPU cores. Each replica
maintains
\begin{equation}
    u(s,a),\qquad v(s,a),\qquad l,
\end{equation}
as well as an empirical state-frequency counter \texttt{state\_freq}.

Within each sampled trajectory, the code performs three coupled updates. The left-eigenvector
update is essentially the same as in Figure 7. The eigenvalue update is also similar. The new
element is the right-eigenvector update:
\begin{equation}
    v(s',a')
    \leftarrow
    v(s',a') + \alpha\left(
        \frac{e^{\beta r}}{l}v(s,a)\times \text{occupancy correction}
        - v(s',a')
    \right),
\end{equation}
where the occupancy correction is built from the prior-policy factors and empirical state
frequencies. The intention is to mimic the forward weighting associated with the right Perron
mode.

This routine does not primarily evaluate the policy itself. Instead it tracks the learned
estimate of
\begin{equation}
    \theta = -\frac{\log l}{\beta},
\end{equation}
stores it periodically, and compares it to the ground-truth model-based value obtained from
\texttt{solve\_unconstrained\_v1}. The final figure has two panels. One panel shows the
learning-rate schedules \(\alpha\) and \(\alpha_\theta\) on a logarithmic scale. The other
panel shows the learned \(\theta\) as a function of training step, averaged over replicas with
a shaded standard-deviation band, together with a dashed horizontal line at the exact target
value. So Figure 8 is a convergence test for sample-based estimation of the large-deviation
rate itself.

Figure 8 uses \(N\) in the same operational sense as Figure 7, but now for the joint learning
of \(u\), \(v\), and \(l\). The code fixes \(N=1000\), wraps the environment with
\texttt{TimeLimit(env, N)}, and then streams one-step transitions from episodes of that maximal
length. Thus \(N\) controls how many updates can be extracted from each episode. In addition,
the bookkeeping intervals are defined from \(N\): the save period is
\texttt{save\_period = N * n\_episodes // 1000}, and the rescaling period is a fixed multiple
of that save period. Therefore even though the horizontal axis of the final plot is
\texttt{training step}, the density of recorded points and the frequency of normalization are
both set indirectly by \(N\).

At the bottom of the file, the \texttt{if \_\_name\_\_ == '\_\_main\_\_'} block simply calls
the figure functions in order. So running \texttt{python main.py} is not ``training a
general agent''. It is executing a sequence of carefully chosen numerical demonstrations of
the paper's claims.

\medskip
\texttt{main.py} is the experimental narrative of the project. It shows how the abstract
objects built in \texttt{utils.py} can be validated by exact propagation, dynamic programming,
Monte Carlo simulation, and sample-based stochastic approximation.

\subsection*{A key distinction: analytical solution versus simulation versus learning}

At this point it is important to separate three very different things that the code does,
because the paper uses all three.

First, there is analytical or model-based solution.
Whenever the code calls
\texttt{get\_dynamics\_and\_rewards},
\texttt{get\_mdp\_transition\_matrix},
or
\texttt{solve\_unconstrained},
it is assuming that the full environment model is known exactly.

That means the code has direct access to:
\begin{equation}
    p(s'|s,a), \qquad r(s,a), \qquad P, \qquad \widetilde{P}.
\end{equation}
From there it solves the Perron problem directly by sparse linear algebra. This is not
reinforcement learning in the usual sample-based sense. It is exact model-based computation on
the finite MDP.

So, for example:
\begin{itemize}
    \item Figure 2 is mostly exact matrix propagation;
    \item Figure 3's thermodynamic curves are obtained by exact Perron solves;
    \item Figure 5's comparison uses exact dynamic programming and exact Perron solves;
    \item Figure 6 is also fully model-based.
\end{itemize}

Second, there is simulation in the narrow sense: rollouts under a fixed policy.
The function \texttt{test\_policy} in \texttt{utils.py} does something much simpler. It does
not learn a policy and does not solve an eigenproblem. It just takes an already specified
policy \(\pi(a|s)\), runs trajectories in the actual environment by repeated calls to
\texttt{env.step(action)}, and averages the realized rewards.

Mathematically, this estimates quantities like
\begin{equation}
    \mathbb{E}_{\tau \sim p_\pi(\tau)}[R_\tau]
\end{equation}
by Monte Carlo sampling.

This part is ``simulation'' in the standard statistical sense: one samples trajectories from a
known process to estimate observables empirically.

Therefore, when Figure 3 compares the analytical curve with ``simulation'', the simulation is
not a separate RL algorithm. It is just Monte Carlo evaluation of the policy that was already
obtained analytically.

Third, there is learning from samples without using the model.
Figures 7 and 8 are the genuinely learning-based part of the repository. There the code does
not build \(\widetilde{P}\) explicitly and does not solve a matrix eigenproblem directly.
Instead, it only sees sampled transitions
\begin{equation}
    (s,a,r,s',a')
\end{equation}
generated by interaction with the environment, and updates tables such as \(u(s,a)\),
\(v(s,a)\), and \(l\) from those samples.

This is much closer to reinforcement learning, because the model is not used explicitly during
the update. However, it is still not ``traditional RL'' in the most common sense:
\begin{itemize}
    \item it is not Q-learning;
    \item it is not policy gradient;
    \item it is not actor-critic;
    \item it is not deep RL.
\end{itemize}

Instead, it is a custom stochastic-approximation algorithm designed to learn the Perron
eigenvectors of the tilted operator from samples.

To make this more concrete, it helps to compare the paper's learning rule with the two most
familiar temporal-difference algorithms.

In Q-learning, the primary object is the action-value function \(Q(s,a)\), and the update tries
to move the current estimate toward
\begin{equation}
    r + \gamma \max_{a'} Q(s',a').
\end{equation}
So Q-learning is a Bellman-backup algorithm, and because the target uses the greedy action at
the next state rather than the action actually taken, it is off-policy.

In SARSA, the primary object is again \(Q(s,a)\), but the update target is
\begin{equation}
    r + \gamma Q(s',a'),
\end{equation}
where \(a'\) is the next action actually sampled by the behavior policy. This is why SARSA is
on-policy. Its name comes from the data tuple
\begin{equation}
    (s,a,r,s',a').
\end{equation}

The paper's \(u\)-\(\theta\) learning shares the \emph{sample format} of SARSA: it also stores
and reuses tuples of the form
\begin{equation}
    (s,a,r,s',a').
\end{equation}
So at the level of replay memory and one-step sampled transitions, it looks closer to SARSA than
to Q-learning.

\subsection{Derivation of Eq.~(25) and the starting point of \(u\)-\(\theta\) learning}

However, the quantity being learned is completely different. The update is not trying to satisfy
a Bellman equation for \(Q\). Instead, it is trying to satisfy the tilted-operator eigenvector
relation
\begin{equation}
    \rho\,u(s,a)
    =
    e^{\beta r(s,a)}
    \mathbb E\!\left[u(s',a')\mid s,a\right],
\end{equation}
with a simultaneous stochastic update for the eigenvalue parameter \(l\) or
\(\theta=-\log l/\beta\). Thus the direct learning targets are
\begin{equation}
    u(s,a), \qquad v(s,a), \qquad \rho \ \text{or}\ \theta,
\end{equation}
not \(Q(s,a)\).

This finite-sample target can be derived directly from the driven kernel formula. Earlier we
obtained
\begin{equation}
    [P_d]_{ji}
    =
    \frac{\widetilde P_{ji}u_j}{e^{-\beta\theta}u_i},
    \qquad
    \widetilde P_{ji}=p(j\mid i)e^{\beta r(i)}.
\end{equation}
Since \(P_d\) is a genuine transition matrix, its columns sum to one:
\begin{equation}
    \sum_j [P_d]_{ji}=1.
\end{equation}
Substituting the previous expression gives
\begin{equation}
    1
    =
    \sum_j
    \frac{p(j\mid i)e^{\beta r(i)}u_j}{e^{-\beta\theta}u_i}
    =
    \frac{e^{\beta r(i)}}{e^{-\beta\theta}u_i}
    \sum_j p(j\mid i)u_j,
\end{equation}
or equivalently
\begin{equation}
    e^{-\beta\theta}u_i
    =
    e^{\beta r(i)}\sum_j p(j\mid i)u_j.
\end{equation}
Writing \(i=(s,a)\) and \(j=(s',a')\), this becomes
\begin{equation}
    e^{-\beta\theta}u(s,a)
    =
    e^{\beta r(s,a)}
    \mathbb E\!\left[u(s',a')\mid s,a\right].
\end{equation}
This is exactly the fixed-point relation used by the \(u\)-\(\theta\) learning rule: one now
replaces the conditional expectation by a single sampled transition \((s,a,r,s',a')\).

For clarity, one full \(u\)-\(\theta\) update can be written as the following workflow:
\begin{enumerate}
    \item sample or replay one tuple
    \begin{equation}
        (s,a,r,s',a');
    \end{equation}
    \item compute the exponential reward weight
    \begin{equation}
        w:=e^{\beta r};
    \end{equation}
    \item read the current table values
    \begin{equation}
        u:=u(s,a),
        \qquad
        u':=u(s',a');
    \end{equation}
    \item form the single-sample fixed-point target
    \begin{equation}
        u_{\mathrm{target}}
        :=
        \frac{w}{l}\,u';
    \end{equation}
    \item update the \(u\)-table entry by exponential averaging
    \begin{equation}
        u(s,a)
        \leftarrow
        (1-\alpha)u(s,a)+\alpha \frac{w}{l}u(s',a');
    \end{equation}
    \item form the single-sample eigenvalue target
    \begin{equation}
        l_{\mathrm{target}}
        :=
        w\,\frac{u'}{u};
    \end{equation}
    \item update the eigenvalue estimate with a slower rate
    \begin{equation}
        l
        \leftarrow
        (1-\alpha_\theta)l+\alpha_\theta w\frac{u'}{u};
    \end{equation}
    \item reconstruct the current policy when needed through
    \begin{equation}
        \pi(a\mid s)\propto \pi_0(a\mid s)u(s,a).
    \end{equation}
\end{enumerate}
This is why the routine is sample-based but not Bellman-based: it is performing a stochastic
approximation of a Perron fixed-point relation rather than a TD backup for \(Q\).

This leads to a different route to the policy. In many entropy-regularized Bellman formulations,
one first learns \(Q\) and then obtains a policy from that value function. In this paper, one
learns the Perron objects first and then reconstructs the policy through
\begin{equation}
    \pi^*(a\mid s)\propto \pi(a\mid s)u(s,a).
\end{equation}
So the pipeline is not
\begin{equation}
    (s,a,r,s',a') \to Q \to \pi,
\end{equation}
but rather
\begin{equation}
    (s,a,r,s',a') \to (u,v,\rho) \to \pi.
\end{equation}

The cleanest summary is:
\begin{itemize}
    \item Q-learning and SARSA live in the Bellman-value-function world;
    \item the paper's \(u\)-\(\theta\) learning lives in the tilted-operator Perron-eigenvector
    world;
    \item the connection to SARSA is in the sampled one-step tuple structure, not in the final
    mathematical target of the update.
\end{itemize}

Traditional reinforcement learning usually tries to learn one of the following:
\begin{equation}
    V^\pi(s), \qquad Q^\pi(s,a), \qquad \pi(a|s),
\end{equation}
or a parameterized approximation of them.

By contrast, this paper's sample-based routines try to learn
\begin{equation}
    u(s,a), \qquad v(s,a), \qquad \rho,
\end{equation}
where \(u\) and \(v\) are the left and right Perron eigenvectors of the tilted operator
\(\widetilde{P}\).

So the learning target is different. The philosophy is:
\begin{equation}
    \text{do not learn } Q \text{ first and then derive the control;}
    \qquad
    \text{learn the tilted-operator eigenstructure directly.}
\end{equation}

The whole project therefore has a very specific logic:
\begin{enumerate}
    \item first solve the finite problem analytically when the model is known;
    \item then simulate trajectories under the analytically derived policy to check that the
    observable averages are correct;
    \item finally, design sample-based updates that recover the same Perron objects without
    explicit access to the model.
\end{enumerate}

This is why the repository can look confusing at first: some figures are exact numerical
linear algebra, some are plain Monte Carlo evaluation, and only some are genuine learning
algorithms. The paper deliberately mixes these because it wants to show both:
\begin{itemize}
    \item the analytical theory is correct;
    \item the same objects can in principle be learned from trajectories.
\end{itemize}

\subsection*{Is this just an interpretation of RL, or a new RL scheme?}

The correct answer is: \emph{both, but not in the same sense}.

At the first level, it is a new theoretical formulation of entropy-regularized control.
At the theoretical level, the paper is not merely re-explaining standard RL notation. It
proposes a different mathematical formulation of the control problem:
\begin{equation}
    \text{control}
    \;\longleftrightarrow\;
    \text{inference on trajectories}
    \;\longleftrightarrow\;
    \text{Perron problem for a tilted operator}.
\end{equation}
So the paper is saying that the optimization target can be understood through biased path
ensembles, large deviations, and dominant eigenvectors rather than only through Bellman
equations.

At the second level, the model-based part is mainly for understanding and exact computation.
When the code explicitly builds \(P\), \(\widetilde{P}\), and solves for \((\rho,u,v)\), this
is not meant to imitate the temporal process of a standard RL agent learning online. It is a
model-based calculation whose main role is:
\begin{itemize}
    \item to make the theory explicit on a finite example;
    \item to check the analytical claims numerically;
    \item to provide a ground truth for comparison.
\end{itemize}
So this part is closer to ``understanding the structure of the control problem exactly'' than
to ``simulating how a conventional RL algorithm learns''.

At the third level, the sample-based part is indeed proposing a different learning route.
Once the paper writes the control problem as an eigenvector problem, it becomes natural to ask:
\emph{can we learn \(u\), \(v\), and \(\rho\) directly from trajectories?}

That is exactly what the Figure 7 and Figure 8 routines try to do. They do not learn
\(Q\)-values first and then derive a policy. Instead they try to learn the tilted-operator
eigenstructure itself:
\begin{equation}
    (u,v,\rho)
    \quad \text{instead of} \quad
    (Q,V,\pi)
    \text{ as the primary object.}
\end{equation}
In that sense, yes, the paper is pointing toward a new reinforcement-learning scheme, or at
least a new family of RL algorithms.

It is not claiming that all of reinforcement learning should be replaced by exact matrix
diagonalization. That would only work when the full finite model is explicitly available.

Instead, the intended logic is:
\begin{enumerate}
    \item derive a new control formulation analytically;
    \item verify it in model-based finite environments;
    \item then design sample-based algorithms inspired by that formulation.
\end{enumerate}

The code is not primarily trying to reproduce the usual RL learning process in order to
explain it. Rather, it develops a different formulation of entropy-regularized RL, solves that
formulation exactly when possible, and then uses those exact solutions to motivate and test a
new eigenvector-based learning approach.

\subsection*{File 4: \texttt{visualization.py} converts mathematical arrays back into spatial pictures}

\paragraph{What this file is for.}
This file contains almost no theory. Its role is representational: it takes arrays living in
state space or state-action space and displays them on the original lattice.

That separation is good design. It ensures that plotting choices do not get mixed into the
numerical or conceptual core.

\paragraph{\texttt{plot\_dist}: one entry point, several display modes.}
The main function \texttt{plot\_dist(...)} is flexible because it accepts different kinds of
objects:
\begin{itemize}
    \item a state distribution over the lattice;
    \item a policy table of shape \((n_S,n_A)\);
    \item a collection of such objects to arrange in multiple subplots;
    \item in some cases a force-field-like object.
\end{itemize}
So the same plotting function can visualize several mathematically different outputs of the
solver.

\paragraph{How a scalar state distribution is displayed.}
If the input can be interpreted as one number per state, the function reshapes it onto the
grid and rescales it between \(0\) and \(1\). Blue intensity is then used to indicate larger
occupancy. This is why the state-distribution plots look like heat maps over the maze.

\paragraph{How a policy is displayed.}
If the input has shape \((n_S,n_A)\), the file interprets it as a policy. Each action is
assigned a direction, and the code draws an arrow from each cell:
\begin{equation}
    \text{arrow length} \propto \pi(a|s).
\end{equation}
Thus the visualization is not a separate optimization object. It is a geometric rendering of
the probability table already computed elsewhere.

\paragraph{How the map semantics are preserved in the figure.}
The helper \texttt{add\_layout} colors walls black. The plotting routines also overlay start,
goal, hole, candy, and nail symbols in different markers and colors. This is more than
cosmetic: without those overlays, one could not visually compare the learned distribution or
policy against the reward and obstacle structure that produced it.

\paragraph{Why this file matters despite being simple.}
The analytical objects of the paper live in high-dimensional indexed arrays. A human, however,
understands the experiment as motion on a maze. \texttt{visualization.py} is the translation
layer between those two viewpoints.

\subsection*{How the four files fit together conceptually}

It is useful to say in one sentence what each file contributes to the derivation:
\begin{itemize}
    \item \texttt{frozen\_lake\_env.py} defines the uncontrolled world;
    \item \texttt{utils.py} rewrites that world as the tilted operator problem of the paper;
    \item \texttt{main.py} asks concrete numerical questions about that operator problem;
    \item \texttt{visualization.py} displays the answers in the original spatial language of the
    maze.
\end{itemize}

Equivalently, the project can be summarized by the following chain:
\begin{equation}
    \texttt{ModifiedFrozenLake}
    \;\Rightarrow\;
    \texttt{env.P}
    \;\Rightarrow\;
    (\texttt{dynamics},\texttt{rewards})
    \;\Rightarrow\;
    P,\widetilde{P}
    \;\Rightarrow\;
    (\rho,u,v)
    \;\Rightarrow\;
    \pi^*,\; \text{driven dynamics},\; \text{bulk distribution}
    \;\Rightarrow\;
    \text{plots and figure comparisons}.
\end{equation}

This is why the code feels unusually close to the mathematics. It is not a broad RL software
framework. It is a narrow computational implementation of one specific analytical story:
entropy-regularized control is rewritten as a Perron-eigenvalue problem for a tilted
state-action transition operator.

\section{Reading Guide: Why NESM and Large Deviations Appear Here}

\subsection*{What NESM means in this paper}

NESM stands for \emph{nonequilibrium statistical mechanics}. The relevant idea here is not
the equilibrium question ``what is the probability of a state?'', but the dynamical question
``what is the probability of a long trajectory?''.

This is exactly why the paper fits naturally into the NESM language. In reinforcement
learning, one studies trajectories
\[
\tau=(z_1,\dots,z_T), \qquad z_t=(s_t,a_t),
\]
and wants to favor those trajectories with large accumulated reward
\[
R_\tau=\sum_{t=1}^T r(z_t).
\]
The control-as-inference reformulation replaces direct reward maximization by an
exponentially tilted trajectory ensemble,
\begin{equation}
    p(\tau\mid O_{1:T})
    \propto
    p(\tau)e^{\beta R_\tau}.
\end{equation}
From the point of view of NESM, this is a standard biased ensemble of trajectories. The
trajectory is reweighted by an exponential of a time-additive observable. That is precisely
the setting where large-deviation theory applies.

\subsection*{Why large deviations appear}

Large-deviation theory studies the probabilities of atypical values of long-time averaged
observables. If
\[
    A_T=\frac{1}{T}\sum_{t=1}^T a(z_t)
\]
is an empirical time average, then large-deviation theory asks how
\[
    P(A_T\approx a)
\]
decays as \(T\to\infty\). The generic form is
\begin{equation}
    P(A_T\approx a)\asymp e^{-T I(a)},
\end{equation}
where \(I(a)\) is the rate function.

In this paper, the additive observable is the total reward or total cost. Instead of fixing
an exact rare value of \(R_\tau/T\), the paper introduces an exponential bias
\[
e^{\beta R_\tau}.
\]
This is the canonical-ensemble version of a large-deviation problem. Therefore:
\begin{itemize}
    \item the tilted matrix \(\widetilde{P}\) is the tilted operator of large-deviation theory;
    \item its dominant eigenvalue plays the role of a scaled cumulant generating function;
    \item the Doob \(h\)-transform gives the driven process that makes the rare trajectories
    typical.
\end{itemize}

So the connection is structural, not metaphorical.

\subsection*{Where large-deviation theory enters the RL derivation}

It is useful to say this more concretely, because otherwise ``large deviation theory'' can
sound like a general slogan rather than a precise ingredient of the derivation. In this paper,
large-deviation theory enters at four specific points.

First, the object of interest is no longer only a policy or a value function, but a
\emph{trajectory ensemble}. After introducing the optimality variables, the paper rewrites the
posterior trajectory weight as
\begin{equation}
    p(\tau\mid O_{1:T})
    \propto
    p(\tau)e^{\beta R_\tau},
\end{equation}
with \(R_\tau=\sum_{t=1}^T r(z_t)\). This is exactly the canonical exponential tilting used in
large-deviation theory for time-additive observables.

Second, the relevant long-time object is the exponential growth rate of the reward-generating
partition function
\begin{equation}
    Z_T(\beta)
    =
    \mathbb E_{p(\tau)}\!\left[e^{\beta R_\tau}\right].
\end{equation}
Large-deviation theory says that for long times one expects
\begin{equation}
    Z_T(\beta)\asymp e^{T\psi(\beta)},
\end{equation}
where \(\psi(\beta)\) is the scaled cumulant generating function. In the paper, the same role is
played by the dominant eigenvalue \(\rho\) of the tilted operator:
\begin{equation}
    \psi(\beta)=\log \rho,
    \qquad
    \theta=-\frac{1}{\beta}\log \rho.
\end{equation}
So the free-energy-like quantity \(\theta\) is not an extra analogy added afterwards; it is the
large-deviation growth rate written in the notation of the paper.

Third, the tilted matrix
\begin{equation}
    \widetilde P_{ji}=P_{ji}e^{\beta r_i}
\end{equation}
is the standard large-deviation tilted operator for a Markov chain with additive observable
\(\sum_t r(z_t)\). The Perron eigenvalue and Perron eigenvectors of \(\widetilde P\) therefore
control the long-time biased trajectory ensemble. This is the Markov-chain version of the
general large-deviation statement that long-time canonical path ensembles are governed by the
dominant spectral data of the tilted generator or tilted transition operator.

Fourth, the driven process is the large-deviation answer to the question: \emph{if the rare
high-reward trajectories are conditioned to occur, what effective Markov dynamics makes them
typical?} The Doob transform gives precisely that effective dynamics. In this sense, the paper's
``optimal controlled dynamics'' is the large-deviation driven process associated with the
reward-biased trajectory ensemble.

So the logic is:
\begin{equation}
    p(\tau)
    \;\longrightarrow\;
    p(\tau)e^{\beta R_\tau}
    \;\longrightarrow\;
    \widetilde P
    \;\longrightarrow\;
    (\rho,u,v)
    \;\longrightarrow\;
    P_d,\;\pi^*.
\end{equation}
Every arrow in this chain is standard large-deviation structure specialized to an
entropy-regularized RL problem.

\subsection*{Exact finite-horizon inference versus long-time large-deviation structure}

It is important not to mix two logically different statements.

The first statement is the control-as-inference equivalence itself. Once the paper defines the
optimality variables by
\begin{equation}
    p(O_t=1\mid z_t)\propto e^{\beta r(z_t)},
\end{equation}
Bayes' rule gives, for any finite horizon \(T\),
\begin{equation}
    p(\tau\mid O_{1:T})
    \propto
    p(\tau)e^{\beta R_\tau}.
\end{equation}
This step is exact. It does not require \(T\to\infty\), Perron-Frobenius theory, driven
processes, or large-deviation asymptotics. It is simply the exact posterior trajectory
distribution induced by the chosen likelihood model for the optimality variables.

The second statement is the long-time structural analysis of that exact posterior. Once one asks
for a time-independent effective dynamics, a bulk stationary distribution, a dominant
exponential growth rate, or equivalence between a constrained and an exponentially tilted
trajectory ensemble, then the long-time limit becomes essential. This is where large-deviation
theory enters:
\begin{itemize}
    \item the tilted operator \(\widetilde P\) controls the long-time posterior weights;
    \item the dominant eigenvalue \(\rho\) gives the asymptotic free-energy or SCGF-like rate;
    \item the eigenvectors \(u,v\) define the bulk distribution and the Doob-transformed driven
    process;
    \item ensemble-equivalence statements belong to this asymptotic layer, not to the basic
    finite-horizon Bayesian identity.
\end{itemize}

So one should read the paper in two stages:
\begin{equation}
    \text{entropy/KL-regularized RL}
    \;\Longrightarrow\;
    \text{exact finite-horizon posterior over trajectories},
\end{equation}
and only afterwards
\begin{equation}
    \text{exact posterior}
    \;+\;
    T\to\infty
    \;\Longrightarrow\;
    \text{large-deviation spectral structure}.
\end{equation}
The first implication is exact and finite-horizon. The second is the additional asymptotic
content that connects the control-as-inference formulation to NESM and large-deviation theory.

\subsection*{Driven process and nonequilibrium steady state}

The phrase \emph{driven process} has a precise meaning here. Start from the uncontrolled Markov
dynamics \(P\), built from the environment transition rule and the prior policy. Under this
original dynamics, trajectories with unusually large total reward are rare. One can still define
them by conditioning on optimality, but that conditioned object is awkward because it is
described through a posterior over whole trajectories.

The point of the Doob transform is to replace this conditioned description by a new Markov
process \(P_d\) such that the trajectories that were rare under \(P\) become typical under
\(P_d\). This new process is the driven process. In RL language, it is the effective controlled
dynamics associated with the optimal policy. In NESM language, it is the auxiliary dynamics that
realizes the biased or conditioned trajectory ensemble without explicitly conditioning at every
step.

This is why the paper writes
\begin{equation}
    [P_d]_{ji}
    =
    \frac{\widetilde P_{ji}u_j}{\rho\,u_i}.
\end{equation}
The original matrix \(P\) is the uncontrolled dynamics. The tilted matrix \(\widetilde P\) is
not yet a true stochastic dynamics because probability is not conserved. The Doob transform uses
the positive left eigenvector \(u\) and the dominant eigenvalue \(\rho\) to convert
\(\widetilde P\) into a genuine stochastic transition matrix \(P_d\). That stochastic matrix is
the driven process.

The phrase \emph{nonequilibrium steady state} refers to the long-time stationary distribution of
this driven dynamics. A steady state means simply that the one-time distribution no longer
changes with time:
\begin{equation}
    p_{t+1}=P_d p_t,
    \qquad
    p_{t+1}=p_t=p_{\mathrm{ss}}.
\end{equation}
But this does \emph{not} imply equilibrium in the sense of detailed balance. The process can
continue to have directional probability currents or cyclic flows while the one-time
distribution remains stationary. That is what makes it a nonequilibrium steady state rather than
an equilibrium state.

In the language of the paper, the bulk distribution
\begin{equation}
    p(z_t\mid O_{1:T})\approx u(z_t)v(z_t)
\end{equation}
is precisely this long-time stationary distribution in the interior of a long trajectory. The
right eigenvector \(v\) describes how surviving or optimal trajectories are distributed forward
in time, while the left eigenvector \(u\) encodes how favorable a present state-action pair is
for future optimality. Their product gives the bulk steady-state weight of the driven process.

So one should distinguish the following three objects carefully:
\begin{itemize}
    \item \(P\): the uncontrolled Markov dynamics;
    \item \(\widetilde P\): the tilted operator that reweights trajectories by reward but is not
    itself probability conserving;
    \item \(P_d\): the driven Markov process obtained from the Doob transform, whose stationary
    distribution in the bulk is \(uv\).
\end{itemize}
The paper's main conceptual move is that the optimal-control problem can be understood through
this driven nonequilibrium steady state.

\subsection*{Why the dominant eigenvalue matters physically}

Once trajectories are reweighted by \(e^{\beta R_\tau}\), the corresponding one-step operator
is no longer a stochastic transition matrix. It becomes the tilted matrix
\[
    \widetilde{P}_{ji}=p(z_{t+1}=j\mid z_t=i)e^{\beta r(i)}.
\]
For long times, powers of \(\widetilde{P}\) are dominated by its Perron eigenvalue
\[
\rho=e^{-\beta\theta}.
\]
This is the same mechanism that appears throughout large-deviation theory:
\begin{itemize}
    \item long-time trajectory weights are controlled by one exponential growth or decay rate;
    \item that rate is encoded in the dominant eigenvalue of the tilted operator;
    \item the associated eigenvectors define the effective driven dynamics that realize the
    biased ensemble.
\end{itemize}

In RL language, that driven dynamics becomes the optimal dynamics and the optimal policy.
In NESM language, it is the auxiliary or driven process conditioned on a rare dynamical
event.

\subsection*{How this relates to nonequilibrium dynamics}

The paper is not saying that RL is identical to all of nonequilibrium physics. The precise
connection is narrower:
\begin{itemize}
    \item the environment together with the prior policy defines a Markov dynamics on
    state-action space;
    \item reward-biasing defines a nonequilibrium trajectory ensemble;
    \item long-time conditioned behavior is described by a driven Markov process obtained from
    the Doob transform.
\end{itemize}

This is a standard pattern in modern nonequilibrium dynamics: start from an unbiased
stochastic dynamics, bias trajectories by a time-integrated observable, then derive the
effective dynamics that makes those rare trajectories typical.

\subsection*{What insight the cited papers can give}

Several of the paper's references are useful because they explain one specific piece of this
logic very clearly.

\paragraph{Touchette, ``The large deviation approach to statistical mechanics''.}
This is the best broad conceptual review. It explains why large-deviation theory naturally
produces the same structures as thermodynamics: entropy functions, free-energy-like objects,
Legendre transforms, and equivalence or inequivalence of ensembles.

\paragraph{Jack and Sollich, ``Large deviations and ensembles of trajectories in stochastic models''.}
This is the most directly relevant bridge from statistical mechanics to the present paper. It
shifts the reader from state ensembles to trajectory ensembles, which is exactly the shift
needed to understand why RL trajectories can be studied with large-deviation tools.

\paragraph{Chetrite and Touchette on nonequilibrium canonical and microcanonical trajectory ensembles.}
These papers explain why one can either:
\begin{itemize}
    \item condition on a rare dynamical event directly; or
    \item bias trajectories exponentially by a conjugate parameter.
\end{itemize}
This is the cleanest background for understanding why the paper moves back and forth between
``optimal trajectories'', ``posterior trajectories'', and ``tilted dynamics''.

\paragraph{Giardina, Kurchan, and Peliti, ``Direct evaluation of large-deviation functions''.}
This paper is useful because it makes large deviations computational. It shows that dominant
eigenvalues and biased trajectory ensembles are not just abstract objects; they can be
computed numerically. This gives good intuition for why the present RL paper turns
eigenvector estimation into a learning problem.

\paragraph{Evans, ``Rules for transition rates in nonequilibrium steady states''.}
This is helpful for intuition about why a ``driven'' or ``auxiliary'' dynamics should exist at
all. If one asks a system to realize an atypical flux or reward structure, the effective
transition rules are modified in a systematic way. That is philosophically close to the role
of the optimal dynamics in the RL paper.

\subsection*{Recommended reading order}

If your goal is specifically to understand this RL paper, the most efficient reading order is:
\begin{enumerate}
    \item Touchette's review, to learn the general large-deviation language;
    \item Jack--Sollich, to understand trajectory ensembles instead of state ensembles;
    \item Chetrite--Touchette, to understand conditioned processes, canonical bias, and Doob
    transforms;
    \item then return to the RL paper and reinterpret \(\widetilde{P}\), \(\rho\), \(u\), \(v\),
    and \(P_d\) in that language.
\end{enumerate}

\subsection*{A one-sentence summary}

The paper uses NESM and large-deviation theory because entropy-regularized RL can be
rewritten as a biased ensemble of trajectories, and the long-time statistics of such biased
ensembles are governed by tilted operators, dominant eigenvalues, and Doob-transformed
dynamics.

\subsection*{What a tensor network is, at the most basic level}

The motivation for tensor networks is simple: a many-body object may formally be one very
large vector, but if its entries are organized by local degrees of freedom, that vector can be
reshaped into a high-order tensor and then approximated by a structured factorization.

Suppose a global configuration is
\begin{equation}
    x=(x_1,x_2,\dots,x_N),
\end{equation}
where each local variable \(x_i\) takes \(d\) values. A scalar function on the full
configuration space,
\begin{equation}
    \psi(x)=\psi(x_1,\dots,x_N),
\end{equation}
can be viewed as a rank-\(N\) tensor
\begin{equation}
    \psi_{x_1x_2\cdots x_N}.
\end{equation}
Writing \(\psi\) as an ordinary vector hides this tensor-product structure:
\begin{equation}
    \mathcal H = \bigotimes_{i=1}^{N} \mathbb{R}^{d}.
\end{equation}

The brute-force representation stores \(d^N\) numbers, which is impossible in the
thermodynamic limit. A tensor network replaces this full table by contracted low-rank local
tensors. In one spatial dimension the standard example is a matrix-product state (MPS):
\begin{equation}
    \psi(x_1,\dots,x_N)
    =
    \sum_{\alpha_1,\dots,\alpha_{N-1}}
    A^{[1]}_{x_1,\alpha_1}
    A^{[2]}_{\alpha_1,x_2,\alpha_2}
    \cdots
    A^{[N]}_{\alpha_{N-1},x_N}.
\end{equation}
The auxiliary indices \(\alpha_i\in\{1,\dots,\chi\}\) are called bond indices, and \(\chi\) is
the bond dimension. If a good approximation exists at moderate \(\chi\), then the number of
parameters grows like \(O(Nd\chi^2)\) instead of \(O(d^N)\).

The same idea applies to operators. A linear map \(K(x',x)\) on the full configuration space
is formally a huge matrix, but if it is built from local rules it can often be represented as a
matrix-product operator (MPO) or, in higher dimensions, by a more general tensor network.

So the basic philosophy is:
\begin{itemize}
    \item a global vector is reinterpreted as a tensor with one physical index per site;
    \item a global matrix is reinterpreted as a tensor with incoming and outgoing physical
    indices per site;
    \item low-rank factorizations are used to exploit locality and weak long-range structure.
\end{itemize}

\subsection*{Why this is relevant here}

In the maze examples of the paper, a state \(s\) is just one lattice position, so the number of
states grows only like \(L^2\). In that setting it is perfectly reasonable to store
\(u(s,a)\), \(v(s,a)\), and \(\widetilde P\) explicitly.

But in a genuine lattice many-body system, a state is usually a full configuration,
\begin{equation}
    s=(s_1,\dots,s_N),
\end{equation}
for example an Ising spin configuration or an occupation-number configuration. Then the state
space size is exponential in \(N\), and the paper's objects become exponentially large:
\begin{equation}
    u(s,a),\qquad v(s,a),\qquad \widetilde P\big((s',a'),(s,a)\big).
\end{equation}
At that point the tabular implementation is no longer meaningful.

This is the context in which one talks about a thermodynamic limit. The question is no longer
``can we enlarge the maze from \(9\times 9\) to \(20\times 20\)?'' but rather:
\begin{equation}
    N\to\infty \quad \text{or} \quad L\to\infty,
\end{equation}
with a local state-update rule and an extensive observable. One then asks whether the dominant
eigenvalue, the driven process, and the optimal policy admit a controlled representation whose
cost grows only polynomially with system size.

\subsection*{How the paper's vectors become tensors}

The simplest case is when there is no explicit action variable and the controlled object is a
Markov process on configurations. Then the right Perron vector is
\begin{equation}
    v(s)=v(s_1,\dots,s_N),
\end{equation}
which can be regarded as a rank-\(N\) tensor
\begin{equation}
    v_{s_1\cdots s_N}.
\end{equation}
Likewise,
\begin{equation}
    u(s)=u(s_1,\dots,s_N)
    \quad \leftrightarrow \quad
    u_{s_1\cdots s_N}.
\end{equation}

If actions are also local, one may fold them into the local physical index and write
\begin{equation}
    z_i=(s_i,a_i),
\end{equation}
so that
\begin{equation}
    v(z)=v(z_1,\dots,z_N)
    \quad \text{and} \quad
    u(z)=u(z_1,\dots,z_N).
\end{equation}
If there is instead one global action coupled to the full configuration, the tensor-product
structure is less natural. So the tensor-network route is most natural when both dynamics and
control are local or quasi-local.

In one dimension, a direct ansatz would be
\begin{equation}
    v(s_1,\dots,s_N)
    \approx
    \sum_{\alpha_1,\dots,\alpha_{N-1}}
    B^{[1]}_{s_1,\alpha_1}
    B^{[2]}_{\alpha_1,s_2,\alpha_2}
    \cdots
    B^{[N]}_{\alpha_{N-1},s_N},
\end{equation}
and similarly for \(u\). In higher dimensions one would need PEPS-like structures, tensor
renormalization, or some other compressed ansatz rather than a simple matrix-product form.

\subsection*{How the tilted operator becomes a tensor-network operator}

The key eigenproblem in the paper is
\begin{equation}
    \widetilde P v = \rho v,
    \qquad
    u^\top \widetilde P = \rho u^\top,
\end{equation}
with
\begin{equation}
    \widetilde P_{ji}=P_{ji}e^{\beta r_i}.
\end{equation}

If the uncontrolled dynamics is local, then \(P\) is not an arbitrary dense matrix. It is built
from local update rules. Likewise, if the reward is additive over sites or over local
transitions,
\begin{equation}
    r(s,a)=\sum_{\ell} r_\ell(\text{local variables}),
\end{equation}
then the factor \(e^{\beta r(s,a)}\) also factorizes locally or quasi-locally. In that case,
\(\widetilde P\) can often be represented as an MPO or another local tensor-network operator.

This is exactly analogous to the transfer-matrix logic of statistical mechanics. The operator
\(\widetilde P\) plays the role of a tilted transfer operator, and \(u\), \(v\), \(\rho\) are
its dominant spectral data. The thermodynamic objects in the paper then become:
\begin{equation}
    \rho \;\; \text{dominant eigenvalue},\qquad
    v \;\; \text{right eigen-tensor},\qquad
    u \;\; \text{left eigen-tensor}.
\end{equation}

So at a formal level, the generalization is very clean:
\begin{equation}
    \text{vector} \to \text{tensor-network state},
    \qquad
    \text{matrix} \to \text{tensor-network operator}.
\end{equation}

\subsection*{How the driven process would be represented}

In the finite-dimensional discussion of the paper, the driven process is obtained by a Doob
transform:
\begin{equation}
    P_d(j\mid i)=\frac{1}{\rho}\frac{v_j}{v_i}\widetilde P_{ji}.
\end{equation}

If \(v_i\) is replaced by a tensor-network approximation \(v_\theta(s)\), then the same formula
still defines a formal driven kernel:
\begin{equation}
    P_d(s'\mid s)=\frac{1}{\rho}\frac{v_\theta(s')}{v_\theta(s)}\widetilde P(s'\mid s).
\end{equation}
So the driven process itself does not disappear in the thermodynamic limit; what changes is the
representation of the eigenfunction \(v\) needed to define it.

The main practical difficulty is that, even if \(v_\theta\) is compact, the ratio
\(v_\theta(s')/v_\theta(s)\) must be computable efficiently for local moves. This is plausible
for local updates and a local tensor ansatz, but it is not automatic. One must choose the
parameterization so that local ratios, local expectations, or effective transition weights can
be evaluated without summing over the full configuration space.

\subsection*{How a tensor-network version of \(u\)-\(\theta\) learning might look}

The paper's current \(u\)-\(\theta\) learning is tabular stochastic approximation. The index
\(\theta\) does not yet denote a neural network or tensor network in the released code. In a
many-body extension, one could let \(\theta\) denote the local tensors of an MPS/MPO-like
ansatz.

At a schematic level, the goal would still be to solve the dominant-eigenvalue problem of the
tilted operator. One would choose parameterized objects
\begin{equation}
    u_\theta,\qquad v_\phi,
\end{equation}
represented by tensor networks, and then minimize a residual or perform a variational power
method so that
\begin{equation}
    \widetilde P v_\phi \approx \rho v_\phi,
    \qquad
    u_\theta^\top \widetilde P \approx \rho u_\theta^\top.
\end{equation}

There are several algorithmic possibilities:
\begin{itemize}
    \item a deterministic variational eigensolver for the tilted MPO;
    \item a power method with repeated compression back into a low-bond-dimension manifold;
    \item a sample-based stochastic update that estimates local eigenvector residuals and
    optimizes the tensor parameters.
\end{itemize}

The third option is the closest in spirit to the paper's current learning viewpoint. But it is
no longer a simple table update. It becomes a variational learning problem on a compressed
manifold of many-body functions.

\subsection*{What is natural, and what is still nontrivial}

Some parts of the generalization are conceptually straightforward.

First, the large-deviation and control-as-inference logic survives unchanged. The exact
finite-horizon posterior identity and the long-time tilted-operator picture do not rely on the
state space being small.

Second, the Perron problem survives unchanged:
\begin{equation}
    \widetilde P v = \rho v.
\end{equation}
Only the representation of \(v\) changes.

Third, the driven process also survives unchanged in formal structure because it is always a
Doob transform of the dominant eigenfunction.

What becomes nontrivial is the approximation theory.

\begin{itemize}
    \item Does \(u\) or \(v\) admit an accurate low-bond-dimension approximation?
    \item Does the tilted operator preserve a manageable local structure after reward biasing?
    \item Can local transition ratios for the Doob transform be computed efficiently?
    \item In two dimensions, is an MPS too restrictive, so that one must move to PEPS-like
    ans\"atze or other compressed representations?
    \item If the reward couples distant sites or imposes a global constraint, does the bond
    dimension have to grow rapidly with system size?
\end{itemize}

So one should not say that tensor networks automatically solve the thermodynamic-limit problem.
What is true is more precise: they provide a plausible representation class for the dominant
eigenvectors and operators when the underlying dynamics and observable are sufficiently local.

\subsection*{A precise interpretation of the supervisor's suggestion}

When your supervisor says that the vectors can be generalized to tensors and treated with
tensor-network ideas, the mathematically precise interpretation is probably the following.

Instead of indexing \(u\) and \(v\) by one global state label,
\begin{equation}
    u_i,\qquad v_i,
\end{equation}
rewrite the state as a product configuration
\begin{equation}
    i \leftrightarrow (s_1,\dots,s_N),
\end{equation}
and reinterpret
\begin{equation}
    u_i \to u_{s_1\cdots s_N},
    \qquad
    v_i \to v_{s_1\cdots s_N}.
\end{equation}
Then replace explicit storage of all components by a compressed factorization, and solve the
tilted Perron problem variationally in that compressed class. The eventual hope is to obtain:
\begin{itemize}
    \item a scalable approximation of the dominant eigenvalue density or free-energy density;
    \item a scalable approximation of the bulk controlled or driven state distribution;
    \item a scalable approximation of local observables under the optimal or driven dynamics.
\end{itemize}

That is a real research program. It is not implemented in the present code, but it is a
mathematically coherent direction for extending the paper from small discrete environments to
many-body thermodynamic limits.

\end{document}
